%%%%%%%%%%%%%%%%%%%%%%%%%%%%%%%%%%%%%%%%%%%%%%%%%%%%%%%%%%%%
% 12.2 AXIOMATIC SET THEORY
% The Composition of Advanced Mathematics
%%%%%%%%%%%%%%%%%%%%%%%%%%%%%%%%%%%%%%%%%%%%%%%%%%%%%%%%%%%%

\section{Axiomatic Set Theory}
\label{sec:axiomatic-set-theory}

\prerequisite{
Elementary set theory, relations, functions, cardinality, infinity,
mathematical induction, formal logic, first-order languages, models,
theories, and methods of proof.
}

\learningobjectives{
By the end of this section, the reader should be able to:
\begin{itemize}
    \item explain why naive set theory leads to paradoxes;
    \item distinguish naive and axiomatic set theory;
    \item understand membership as the primitive relation of set theory;
    \item recognize the language of set theory;
    \item understand extensionality;
    \item understand the empty-set axiom;
    \item understand pairing;
    \item understand union;
    \item understand power sets;
    \item understand separation schemas;
    \item understand replacement schemas;
    \item understand the axiom of infinity;
    \item understand regularity or foundation;
    \item understand the axiom of choice;
    \item recognize the Zermelo--Fraenkel axioms;
    \item distinguish ZF from ZFC;
    \item understand why unrestricted comprehension fails;
    \item derive basic set constructions from axioms;
    \item define ordered pairs set-theoretically;
    \item define Cartesian products;
    \item define relations as sets;
    \item define functions as special relations;
    \item understand set-theoretic representations of numbers;
    \item construct the von Neumann natural numbers;
    \item understand finite ordinals;
    \item define successor ordinals;
    \item understand transitive sets;
    \item recognize ordinals;
    \item understand ordinal ordering;
    \item distinguish ordinal and cardinal concepts;
    \item recognize transfinite induction;
    \item understand transfinite recursion conceptually;
    \item define equipotence;
    \item understand cardinal numbers conceptually;
    \item distinguish finite, countable, and uncountable sets;
    \item understand Cantor's theorem;
    \item understand diagonal arguments;
    \item recognize cardinal exponentiation;
    \item understand the continuum;
    \item recognize the continuum hypothesis;
    \item distinguish theorem, axiom, and independence statement;
    \item understand well-ordering;
    \item recognize equivalent forms of the axiom of choice;
    \item understand Zorn's lemma;
    \item understand the well-ordering theorem;
    \item recognize maximality principles;
    \item understand filters and ultrafilters conceptually;
    \item recognize transitive closure;
    \item understand rank and the cumulative hierarchy;
    \item recognize the sets \(V_\alpha\);
    \item understand why the cumulative hierarchy avoids self-membership;
    \item understand proper classes conceptually;
    \item distinguish sets from proper classes;
    \item recognize examples such as the class of all sets and all ordinals;
    \item understand relative consistency conceptually;
    \item recognize independence phenomena;
    \item understand the role of models of set theory;
    \item recognize constructible sets conceptually;
    \item understand forcing conceptually;
    \item recognize large cardinal axioms conceptually;
    \item connect axiomatic set theory with logic, algebra, topology,
          analysis, model theory, proof theory, category theory, and the
          foundations of mathematics.
\end{itemize}
}

%%%%%%%%%%%%%%%%%%%%%%%%%%%%%%%%%%%%%%%%%%%%%%%%%%%%%%%%%%%%
\subsection{Why Axiomatize Set Theory?}
\label{subsec:why-axiomatize-set-theory}
%%%%%%%%%%%%%%%%%%%%%%%%%%%%%%%%%%%%%%%%%%%%%%%%%%%%%%%%%%%%

Elementary mathematics often treats a set as an arbitrary collection of
objects.

This informal viewpoint is useful, but unrestricted collection formation
leads to contradictions.

Axiomatic set theory replaces the statement

\[
\boxed{
\text{``any describable collection is a set''}
}
\]

with a controlled system of axioms specifying which sets may be formed.

%%%%%%%%%%%%%%%%%%%%%%%%%%%%%%%%%%%%%%%%%%%%%%%%%%%%%%%%%%%%
\subsection{Naive Set Theory}
\label{subsec:naive-set-theory}
%%%%%%%%%%%%%%%%%%%%%%%%%%%%%%%%%%%%%%%%%%%%%%%%%%%%%%%%%%%%

Naive set theory informally permits constructions such as

\[
\boxed{
\{x:P(x)\},
}
\]

for an arbitrary property \(P\).

This is called unrestricted comprehension.

It appears natural, but it is inconsistent.

The failure is demonstrated by Russell's paradox.

%%%%%%%%%%%%%%%%%%%%%%%%%%%%%%%%%%%%%%%%%%%%%%%%%%%%%%%%%%%%
\subsection{Russell's Paradox}
\label{subsec:russells-paradox}
%%%%%%%%%%%%%%%%%%%%%%%%%%%%%%%%%%%%%%%%%%%%%%%%%%%%%%%%%%%%

Consider the supposed set

\[
\boxed{
R
=
\{x:x\notin x\}.
}
\]

Ask whether

\[
R\in R.
\]

By definition,

\[
R\in R
\iff
R\notin R.
\]

Thus

\[
\boxed{
R\in R
\iff
R\notin R,
}
\]

which is contradictory.

Therefore unrestricted comprehension cannot be allowed.

%%%%%%%%%%%%%%%%%%%%%%%%%%%%%%%%%%%%%%%%%%%%%%%%%%%%%%%%%%%%
\subsection{The Axiomatic Response}
\label{subsec:axiomatic-response}
%%%%%%%%%%%%%%%%%%%%%%%%%%%%%%%%%%%%%%%%%%%%%%%%%%%%%%%%%%%%

Modern set theory begins with a formal language and a collection of axioms.

The standard theory used throughout most mathematics is

\[
\boxed{
\mathrm{ZFC},
}
\]

which stands for

\[
\boxed{
\text{Zermelo--Fraenkel set theory with the Axiom of Choice}.
}
\]

Without Choice, the theory is denoted

\[
\boxed{
\mathrm{ZF}.
}
\]

%%%%%%%%%%%%%%%%%%%%%%%%%%%%%%%%%%%%%%%%%%%%%%%%%%%%%%%%%%%%
\subsection{Language of Set Theory}
\label{subsec:language-of-set-theory}
%%%%%%%%%%%%%%%%%%%%%%%%%%%%%%%%%%%%%%%%%%%%%%%%%%%%%%%%%%%%

The usual first-order language of pure set theory contains one
nonlogical binary relation symbol:

\[
\boxed{
\in.
}
\]

The formula

\[
x\in y
\]

means that \(x\) is an element of \(y\).

Other familiar notation such as

\[
\subseteq,
\qquad
\varnothing,
\qquad
\cup,
\qquad
\mathcal{P},
\]

can be defined from membership and logical notation.

%%%%%%%%%%%%%%%%%%%%%%%%%%%%%%%%%%%%%%%%%%%%%%%%%%%%%%%%%%%%
\subsection{Equality in Set Theory}
\label{subsec:set-theoretic-equality}
%%%%%%%%%%%%%%%%%%%%%%%%%%%%%%%%%%%%%%%%%%%%%%%%%%%%%%%%%%%%

Equality is the ordinary logical equality symbol.

The defining structural principle of sets is that a set is determined
entirely by its elements.

This is formalized by the axiom of extensionality.

%%%%%%%%%%%%%%%%%%%%%%%%%%%%%%%%%%%%%%%%%%%%%%%%%%%%%%%%%%%%
\subsection{Axiom of Extensionality}
\label{subsec:axiom-extensionality}
%%%%%%%%%%%%%%%%%%%%%%%%%%%%%%%%%%%%%%%%%%%%%%%%%%%%%%%%%%%%

\begin{definitionbox}{Axiom of Extensionality}
Two sets are equal exactly when they have the same elements:

\[
\boxed{
\forall A\forall B
\left[
\left(
\forall x
\,
(x\in A\leftrightarrow x\in B)
\right)
\rightarrow
A=B
\right].
}
\]
\end{definitionbox}

This means sets have no hidden internal structure beyond their membership.

%%%%%%%%%%%%%%%%%%%%%%%%%%%%%%%%%%%%%%%%%%%%%%%%%%%%%%%%%%%%
\subsection{Consequences of Extensionality}
\label{subsec:extensionality-consequences}
%%%%%%%%%%%%%%%%%%%%%%%%%%%%%%%%%%%%%%%%%%%%%%%%%%%%%%%%%%%%

If

\[
A=\{1,2,3\}
\]

and

\[
B=\{3,2,1\},
\]

then

\[
\boxed{
A=B.
}
\]

Order does not matter.

Similarly,

\[
\{1,1,2,2,3\}
\]

has the same elements as

\[
\{1,2,3\},
\]

so extensionality gives

\[
\boxed{
\{1,1,2,2,3\}
=
\{1,2,3\}.
}
\]

Multiplicity does not matter.

%%%%%%%%%%%%%%%%%%%%%%%%%%%%%%%%%%%%%%%%%%%%%%%%%%%%%%%%%%%%
\subsection{Empty Set}
\label{subsec:axiomatic-empty-set}
%%%%%%%%%%%%%%%%%%%%%%%%%%%%%%%%%%%%%%%%%%%%%%%%%%%%%%%%%%%%

ZF proves the existence of a set having no elements.

One commonly presents this through an empty-set axiom:

\[
\boxed{
\exists A
\forall x
\,
(x\notin A).
}
\]

By extensionality, any two sets with no elements are equal.

This unique set is denoted

\[
\boxed{
\varnothing.
}
\]

%%%%%%%%%%%%%%%%%%%%%%%%%%%%%%%%%%%%%%%%%%%%%%%%%%%%%%%%%%%%
\subsection{Worked Example: Uniqueness of the Empty Set}
\label{subsec:worked-empty-set-uniqueness}
%%%%%%%%%%%%%%%%%%%%%%%%%%%%%%%%%%%%%%%%%%%%%%%%%%%%%%%%%%%%

\begin{workedexamplebox}{The Empty Set Is Unique}

Suppose \(A\) and \(B\) both have no elements.

Then for every \(x\),

\[
x\in A
\]

is false and

\[
x\in B
\]

is false.

Therefore,

\[
x\in A
\leftrightarrow
x\in B.
\]

By extensionality,

\begin{answerbox}
\[
\boxed{
A=B.
}
\]
\end{answerbox}

Thus the empty set is unique.

\end{workedexamplebox}

%%%%%%%%%%%%%%%%%%%%%%%%%%%%%%%%%%%%%%%%%%%%%%%%%%%%%%%%%%%%
\subsection{Axiom of Pairing}
\label{subsec:axiom-pairing}
%%%%%%%%%%%%%%%%%%%%%%%%%%%%%%%%%%%%%%%%%%%%%%%%%%%%%%%%%%%%

\begin{definitionbox}{Axiom of Pairing}
For any sets \(a\) and \(b\), there exists a set whose elements are exactly
\(a\) and \(b\):

\[
\boxed{
\forall a\forall b
\exists c
\forall x
\left[
x\in c
\leftrightarrow
(x=a\lor x=b)
\right].
}
\]
\end{definitionbox}

The resulting set is denoted

\[
\boxed{
\{a,b\}.
}
\]

%%%%%%%%%%%%%%%%%%%%%%%%%%%%%%%%%%%%%%%%%%%%%%%%%%%%%%%%%%%%
\subsection{Singleton Sets}
\label{subsec:singleton-construction}
%%%%%%%%%%%%%%%%%%%%%%%%%%%%%%%%%%%%%%%%%%%%%%%%%%%%%%%%%%%%

Set

\[
a=b.
\]

Pairing then produces

\[
\boxed{
\{a,a\}
=
\{a\}.
}
\]

Thus singleton sets need no separate axiom.

%%%%%%%%%%%%%%%%%%%%%%%%%%%%%%%%%%%%%%%%%%%%%%%%%%%%%%%%%%%%
\subsection{Axiom of Union}
\label{subsec:axiom-union}
%%%%%%%%%%%%%%%%%%%%%%%%%%%%%%%%%%%%%%%%%%%%%%%%%%%%%%%%%%%%

\begin{definitionbox}{Axiom of Union}
For every set \(A\), there exists a set containing exactly the elements of
the elements of \(A\):

\[
\boxed{
\forall A
\exists U
\forall x
\left[
x\in U
\leftrightarrow
\exists B
(B\in A\land x\in B)
\right].
}
\]
\end{definitionbox}

This set is denoted

\[
\boxed{
\bigcup A.
}
\]

%%%%%%%%%%%%%%%%%%%%%%%%%%%%%%%%%%%%%%%%%%%%%%%%%%%%%%%%%%%%
\subsection{Worked Example: Union}
\label{subsec:worked-axiomatic-union}
%%%%%%%%%%%%%%%%%%%%%%%%%%%%%%%%%%%%%%%%%%%%%%%%%%%%%%%%%%%%

\begin{workedexamplebox}{Union of a Set of Sets}

Let

\[
A
=
\{
\{1,2\},
\{2,3\},
\{4\}
\}.
\]

Then

\[
\bigcup A
\]

contains every element appearing inside one of the members of \(A\).

Therefore,

\begin{answerbox}
\[
\boxed{
\bigcup A
=
\{1,2,3,4\}.
}
\]
\end{answerbox}

\end{workedexamplebox}

%%%%%%%%%%%%%%%%%%%%%%%%%%%%%%%%%%%%%%%%%%%%%%%%%%%%%%%%%%%%
\subsection{Binary Union}
\label{subsec:binary-union-set-theory}
%%%%%%%%%%%%%%%%%%%%%%%%%%%%%%%%%%%%%%%%%%%%%%%%%%%%%%%%%%%%

Using pairing and union, define

\[
\boxed{
A\cup B
=
\bigcup
\{A,B\}.
}
\]

Thus ordinary binary union is derived from more primitive axioms.

%%%%%%%%%%%%%%%%%%%%%%%%%%%%%%%%%%%%%%%%%%%%%%%%%%%%%%%%%%%%
\subsection{Axiom of Power Set}
\label{subsec:axiom-power-set}
%%%%%%%%%%%%%%%%%%%%%%%%%%%%%%%%%%%%%%%%%%%%%%%%%%%%%%%%%%%%

\begin{definitionbox}{Axiom of Power Set}
For every set \(A\), there exists a set containing exactly all subsets of
\(A\):

\[
\boxed{
\forall A
\exists P
\forall x
\left[
x\in P
\leftrightarrow
x\subseteq A
\right].
}
\]
\end{definitionbox}

This set is denoted

\[
\boxed{
\mathcal{P}(A).
}
\]

%%%%%%%%%%%%%%%%%%%%%%%%%%%%%%%%%%%%%%%%%%%%%%%%%%%%%%%%%%%%
\subsection{Subset Relation}
\label{subsec:subset-definition-axiomatic}
%%%%%%%%%%%%%%%%%%%%%%%%%%%%%%%%%%%%%%%%%%%%%%%%%%%%%%%%%%%%

The notation

\[
A\subseteq B
\]

abbreviates

\[
\boxed{
\forall x
\,
(x\in A\rightarrow x\in B).
}
\]

Thus subset is not a primitive relation in the language of pure set theory.

It is defined using membership and logic.

%%%%%%%%%%%%%%%%%%%%%%%%%%%%%%%%%%%%%%%%%%%%%%%%%%%%%%%%%%%%
\subsection{Worked Example: Power Set}
\label{subsec:worked-power-set}
%%%%%%%%%%%%%%%%%%%%%%%%%%%%%%%%%%%%%%%%%%%%%%%%%%%%%%%%%%%%

\begin{workedexamplebox}{Power Set of a Two-Element Set}

Let

\[
A=\{a,b\}.
\]

Its subsets are

\[
\varnothing,
\qquad
\{a\},
\qquad
\{b\},
\qquad
\{a,b\}.
\]

Therefore,

\begin{answerbox}
\[
\boxed{
\mathcal{P}(A)
=
\{
\varnothing,
\{a\},
\{b\},
\{a,b\}
\}.
}
\]
\end{answerbox}

\end{workedexamplebox}

%%%%%%%%%%%%%%%%%%%%%%%%%%%%%%%%%%%%%%%%%%%%%%%%%%%%%%%%%%%%
\subsection{Restricted Comprehension}
\label{subsec:restricted-comprehension}
%%%%%%%%%%%%%%%%%%%%%%%%%%%%%%%%%%%%%%%%%%%%%%%%%%%%%%%%%%%%

Russell's paradox shows that unrestricted comprehension is inconsistent.

Instead, ZF uses restricted set formation.

Given an existing set \(A\), one may select elements satisfying a property:

\[
\boxed{
\{x\in A:\varphi(x)\}.
}
\]

This is governed by the Separation Schema.

%%%%%%%%%%%%%%%%%%%%%%%%%%%%%%%%%%%%%%%%%%%%%%%%%%%%%%%%%%%%
\subsection{Separation Schema}
\label{subsec:separation-schema}
%%%%%%%%%%%%%%%%%%%%%%%%%%%%%%%%%%%%%%%%%%%%%%%%%%%%%%%%%%%%

For each first-order formula \(\varphi(x,\mathbf{p})\), ZF includes an
axiom asserting that

\[
\boxed{
\{x\in A:\varphi(x,\mathbf{p})\}
}
\]

exists.

Formally,

\[
\boxed{
\forall A
\exists B
\forall x
\left[
x\in B
\leftrightarrow
(x\in A\land\varphi(x,\mathbf{p}))
\right].
}
\]

Because this is supplied for every suitable formula \(\varphi\), it is an
axiom schema rather than one individual sentence in the ordinary
presentation.

%%%%%%%%%%%%%%%%%%%%%%%%%%%%%%%%%%%%%%%%%%%%%%%%%%%%%%%%%%%%
\subsection{Why Separation Avoids Russell's Paradox}
\label{subsec:separation-russell}
%%%%%%%%%%%%%%%%%%%%%%%%%%%%%%%%%%%%%%%%%%%%%%%%%%%%%%%%%%%%

Separation does not allow

\[
\{x:x\notin x\}
\]

to be formed from an unrestricted universal domain.

It permits only

\[
\boxed{
\{x\in A:x\notin x\}
}
\]

for an already existing set \(A\).

Thus one cannot automatically assume the existence of a set containing
every set.

%%%%%%%%%%%%%%%%%%%%%%%%%%%%%%%%%%%%%%%%%%%%%%%%%%%%%%%%%%%%
\subsection{Intersection}
\label{subsec:axiomatic-intersection}
%%%%%%%%%%%%%%%%%%%%%%%%%%%%%%%%%%%%%%%%%%%%%%%%%%%%%%%%%%%%

If \(A\neq\varnothing\), define

\[
\boxed{
\bigcap A
=
\{
x\in B:
\forall C\in A,\,
x\in C
\},
}
\]

where \(B\) is any fixed member of \(A\).

Separation produces the desired subset.

For binary intersections,

\[
\boxed{
A\cap B
=
\{
x\in A:x\in B
\}.
}
\]

%%%%%%%%%%%%%%%%%%%%%%%%%%%%%%%%%%%%%%%%%%%%%%%%%%%%%%%%%%%%
\subsection{Set Difference}
\label{subsec:axiomatic-set-difference}
%%%%%%%%%%%%%%%%%%%%%%%%%%%%%%%%%%%%%%%%%%%%%%%%%%%%%%%%%%%%

Set difference is obtained by separation:

\[
\boxed{
A\setminus B
=
\{
x\in A:x\notin B
\}.
}
\]

No separate set-difference axiom is required.

%%%%%%%%%%%%%%%%%%%%%%%%%%%%%%%%%%%%%%%%%%%%%%%%%%%%%%%%%%%%
\subsection{Axiom Schema of Replacement}
\label{subsec:replacement-schema}
%%%%%%%%%%%%%%%%%%%%%%%%%%%%%%%%%%%%%%%%%%%%%%%%%%%%%%%%%%%%

Replacement allows images of sets under definable functional rules to be
collected into sets.

Suppose formula \(\varphi(x,y)\) defines a unique \(y\) for every \(x\in A\).

Then Replacement asserts the existence of

\[
\boxed{
\{
y:
\exists x\in A\,\varphi(x,y)
\}.
}
\]

Schematically,

\[
\boxed{
\text{set-sized domain}
+
\text{definable functional rule}
\Longrightarrow
\text{set-sized image}.
}
\]

%%%%%%%%%%%%%%%%%%%%%%%%%%%%%%%%%%%%%%%%%%%%%%%%%%%%%%%%%%%%
\subsection{Why Replacement Matters}
\label{subsec:why-replacement-matters}
%%%%%%%%%%%%%%%%%%%%%%%%%%%%%%%%%%%%%%%%%%%%%%%%%%%%%%%%%%%%

Replacement is important when iterating constructions over large ordinal
domains.

For a set

\[
A
\]

and definable function

\[
F,
\]

Replacement supports forming

\[
\boxed{
F[A]
=
\{F(x):x\in A\}.
}
\]

It greatly strengthens the ability of ZF to carry out transfinite
constructions.

%%%%%%%%%%%%%%%%%%%%%%%%%%%%%%%%%%%%%%%%%%%%%%%%%%%%%%%%%%%%
\subsection{Separation Versus Replacement}
\label{subsec:separation-vs-replacement}
%%%%%%%%%%%%%%%%%%%%%%%%%%%%%%%%%%%%%%%%%%%%%%%%%%%%%%%%%%%%

Separation constructs a subset of an existing set:

\[
\boxed{
B
\subseteq
A.
}
\]

Replacement takes elements of a set and replaces each by a uniquely
defined output:

\[
\boxed{
x
\mapsto
F(x).
}
\]

The outputs need not already lie inside one known containing set.

%%%%%%%%%%%%%%%%%%%%%%%%%%%%%%%%%%%%%%%%%%%%%%%%%%%%%%%%%%%%
\subsection{Axiom of Infinity}
\label{subsec:axiom-infinity}
%%%%%%%%%%%%%%%%%%%%%%%%%%%%%%%%%%%%%%%%%%%%%%%%%%%%%%%%%%%%

Set theory must explicitly assert that infinite sets exist.

One common form of the Axiom of Infinity says there exists a set \(I\) such
that

\[
\boxed{
\varnothing\in I
}
\]

and whenever

\[
x\in I,
\]

its successor

\[
\boxed{
x\cup\{x\}
}
\]

also belongs to \(I\).

Formally,

\[
\boxed{
\exists I
\left[
\varnothing\in I
\land
\forall x
\left(
x\in I
\rightarrow
x\cup\{x\}\in I
\right)
\right].
}
\]

%%%%%%%%%%%%%%%%%%%%%%%%%%%%%%%%%%%%%%%%%%%%%%%%%%%%%%%%%%%%
\subsection{Inductive Sets}
\label{subsec:inductive-sets}
%%%%%%%%%%%%%%%%%%%%%%%%%%%%%%%%%%%%%%%%%%%%%%%%%%%%%%%%%%%%

\begin{definitionbox}{Inductive Set}
A set \(I\) is inductive if

\[
\varnothing\in I
\]

and

\[
x\in I
\Longrightarrow
x\cup\{x\}\in I.
\]
\end{definitionbox}

The Axiom of Infinity asserts that at least one inductive set exists.

%%%%%%%%%%%%%%%%%%%%%%%%%%%%%%%%%%%%%%%%%%%%%%%%%%%%%%%%%%%%
\subsection{Von Neumann Natural Numbers}
\label{subsec:von-neumann-natural-numbers}
%%%%%%%%%%%%%%%%%%%%%%%%%%%%%%%%%%%%%%%%%%%%%%%%%%%%%%%%%%%%

Define

\[
\boxed{
0
=
\varnothing.
}
\]

For each natural number \(n\), define its successor by

\[
\boxed{
n+1
=
n\cup\{n\}.
}
\]

Then:

\[
0=\varnothing,
\]

\[
1=\{0\},
\]

\[
2=\{0,1\},
\]

\[
3=\{0,1,2\}.
\]

Thus each natural number is the set of all smaller natural numbers.

%%%%%%%%%%%%%%%%%%%%%%%%%%%%%%%%%%%%%%%%%%%%%%%%%%%%%%%%%%%%
\subsection{Worked Example: First Natural Numbers}
\label{subsec:worked-von-neumann-naturals}
%%%%%%%%%%%%%%%%%%%%%%%%%%%%%%%%%%%%%%%%%%%%%%%%%%%%%%%%%%%%

\begin{workedexamplebox}{Building \(0,1,2,3\)}

Start with

\[
0=\varnothing.
\]

Then

\[
1
=
0\cup\{0\}
=
\{\varnothing\}.
\]

Next,

\[
2
=
1\cup\{1\}
=
\{0,1\}.
\]

Then,

\[
3
=
2\cup\{2\}
=
\{0,1,2\}.
\]

\begin{answerbox}
\[
\boxed{
0=\varnothing,
\qquad
1=\{0\},
\qquad
2=\{0,1\},
\qquad
3=\{0,1,2\}.
}
\]
\end{answerbox}

\end{workedexamplebox}

%%%%%%%%%%%%%%%%%%%%%%%%%%%%%%%%%%%%%%%%%%%%%%%%%%%%%%%%%%%%
\subsection{The Set \(\omega\)}
\label{subsec:set-omega}
%%%%%%%%%%%%%%%%%%%%%%%%%%%%%%%%%%%%%%%%%%%%%%%%%%%%%%%%%%%%

The natural numbers are collected into the least inductive set.

This set is denoted

\[
\boxed{
\omega.
}
\]

Conceptually,

\[
\boxed{
\omega
=
\{
0,1,2,3,\ldots
\}.
}
\]

The existence of \(\omega\) is obtained from Infinity together with
Separation.

%%%%%%%%%%%%%%%%%%%%%%%%%%%%%%%%%%%%%%%%%%%%%%%%%%%%%%%%%%%%
\subsection{Finite Ordinals}
\label{subsec:finite-ordinals}
%%%%%%%%%%%%%%%%%%%%%%%%%%%%%%%%%%%%%%%%%%%%%%%%%%%%%%%%%%%%

Under the von Neumann representation,

\[
\boxed{
n
=
\{0,1,\ldots,n-1\}.
}
\]

Hence if

\[
m<n,
\]

then

\[
\boxed{
m\in n.
}
\]

Membership itself represents the usual strict ordering among natural
numbers.

%%%%%%%%%%%%%%%%%%%%%%%%%%%%%%%%%%%%%%%%%%%%%%%%%%%%%%%%%%%%
\subsection{Axiom of Foundation}
\label{subsec:axiom-foundation}
%%%%%%%%%%%%%%%%%%%%%%%%%%%%%%%%%%%%%%%%%%%%%%%%%%%%%%%%%%%%

The Axiom of Foundation, also called Regularity, asserts that every
nonempty set \(A\) contains an element disjoint from \(A\):

\[
\boxed{
\forall A
\left[
A\neq\varnothing
\rightarrow
\exists x
\left(
x\in A
\land
x\cap A=\varnothing
\right)
\right].
}
\]

This rules out certain membership cycles.

%%%%%%%%%%%%%%%%%%%%%%%%%%%%%%%%%%%%%%%%%%%%%%%%%%%%%%%%%%%%
\subsection{No Self-Membership}
\label{subsec:no-self-membership}
%%%%%%%%%%%%%%%%%%%%%%%%%%%%%%%%%%%%%%%%%%%%%%%%%%%%%%%%%%%%

Foundation implies

\[
\boxed{
x\notin x
}
\]

for every set \(x\).

Suppose instead that

\[
x\in x.
\]

Then consider

\[
A=\{x\}.
\]

Its only element \(x\) satisfies

\[
x\cap A\neq\varnothing
\]

because

\[
x\in x
\]

and

\[
x\in A.
\]

This contradicts Foundation.

%%%%%%%%%%%%%%%%%%%%%%%%%%%%%%%%%%%%%%%%%%%%%%%%%%%%%%%%%%%%
\subsection{No Finite Membership Cycles}
\label{subsec:no-membership-cycles}
%%%%%%%%%%%%%%%%%%%%%%%%%%%%%%%%%%%%%%%%%%%%%%%%%%%%%%%%%%%%

Foundation also rules out cycles such as

\[
x_0\in x_1\in\cdots\in x_n\in x_0.
\]

Thus the membership relation is well-founded in the standard set-theoretic
universe.

This motivates the cumulative hierarchy.

%%%%%%%%%%%%%%%%%%%%%%%%%%%%%%%%%%%%%%%%%%%%%%%%%%%%%%%%%%%%
\subsection{Axiom of Choice}
\label{subsec:axiom-choice}
%%%%%%%%%%%%%%%%%%%%%%%%%%%%%%%%%%%%%%%%%%%%%%%%%%%%%%%%%%%%

The Axiom of Choice states that for every set \(X\) of nonempty sets, there
exists a function selecting one element from each member of \(X\).

Formally, if

\[
\forall A\in X,
\qquad
A\neq\varnothing,
\]

then there exists a function \(f\) such that

\[
\boxed{
f(A)\in A
}
\]

for every

\[
A\in X.
\]

Such an \(f\) is called a choice function.

%%%%%%%%%%%%%%%%%%%%%%%%%%%%%%%%%%%%%%%%%%%%%%%%%%%%%%%%%%%%
\subsection{Finite Versus Infinite Choice}
\label{subsec:finite-infinite-choice}
%%%%%%%%%%%%%%%%%%%%%%%%%%%%%%%%%%%%%%%%%%%%%%%%%%%%%%%%%%%%

For finitely many nonempty sets, one can make finitely many selections
without invoking the full Axiom of Choice.

The significance of Choice appears for arbitrary infinite families.

There may be no explicit rule determining which element to select from each
set.

Choice asserts that a global selection nevertheless exists.

%%%%%%%%%%%%%%%%%%%%%%%%%%%%%%%%%%%%%%%%%%%%%%%%%%%%%%%%%%%%
\subsection{Worked Example: Choice Function}
\label{subsec:worked-choice-function}
%%%%%%%%%%%%%%%%%%%%%%%%%%%%%%%%%%%%%%%%%%%%%%%%%%%%%%%%%%%%

\begin{workedexamplebox}{Selecting From a Family}

Suppose

\[
X
=
\{
\{1,2\},
\{a,b,c\},
\{\pi,e\}
\}.
\]

One possible choice function is

\[
f(\{1,2\})=1,
\]

\[
f(\{a,b,c\})=a,
\]

and

\[
f(\{\pi,e\})=\pi.
\]

\begin{answerbox}
\[
\boxed{
f(A)\in A
\quad
\text{for every }A\in X.
}
\]
\end{answerbox}

For finite families this construction is straightforward; the Axiom of
Choice concerns arbitrary families.

\end{workedexamplebox}

%%%%%%%%%%%%%%%%%%%%%%%%%%%%%%%%%%%%%%%%%%%%%%%%%%%%%%%%%%%%
\subsection{Zermelo--Fraenkel Set Theory}
\label{subsec:zf-set-theory}
%%%%%%%%%%%%%%%%%%%%%%%%%%%%%%%%%%%%%%%%%%%%%%%%%%%%%%%%%%%%

The standard ZF axioms are usually organized around:

\[
\boxed{
\text{Extensionality},
}
\]

\[
\boxed{
\text{Pairing},
}
\]

\[
\boxed{
\text{Union},
}
\]

\[
\boxed{
\text{Power Set},
}
\]

\[
\boxed{
\text{Infinity},
}
\]

\[
\boxed{
\text{Separation},
}
\]

\[
\boxed{
\text{Replacement},
}
\]

and

\[
\boxed{
\text{Foundation}.
}
\]

An empty set can be derived in common formulations, though some
presentations state its existence separately for convenience.

%%%%%%%%%%%%%%%%%%%%%%%%%%%%%%%%%%%%%%%%%%%%%%%%%%%%%%%%%%%%
\subsection{ZFC}
\label{subsec:zfc}
%%%%%%%%%%%%%%%%%%%%%%%%%%%%%%%%%%%%%%%%%%%%%%%%%%%%%%%%%%%%

Adding the Axiom of Choice gives

\[
\boxed{
\mathrm{ZFC}
=
\mathrm{ZF}
+
\mathrm{AC}.
}
\]

ZFC is the most widely used foundational framework for ordinary modern
mathematics.

Most standard mathematical constructions can be formalized within it.

%%%%%%%%%%%%%%%%%%%%%%%%%%%%%%%%%%%%%%%%%%%%%%%%%%%%%%%%%%%%
\subsection{Derived Set Constructions}
\label{subsec:derived-set-constructions}
%%%%%%%%%%%%%%%%%%%%%%%%%%%%%%%%%%%%%%%%%%%%%%%%%%%%%%%%%%%%

Many familiar set operations do not require separate axioms.

From the ZF axioms one can define and prove the existence of:

\[
\boxed{
\{a\},
}
\]

\[
\boxed{
A\cup B,
}
\]

\[
\boxed{
A\cap B,
}
\]

\[
\boxed{
A\setminus B,
}
\]

and various Cartesian products and function spaces.

The axioms provide a small foundation from which more elaborate
constructions are derived.

%%%%%%%%%%%%%%%%%%%%%%%%%%%%%%%%%%%%%%%%%%%%%%%%%%%%%%%%%%%%
\subsection{Ordered Pairs}
\label{subsec:set-theoretic-ordered-pairs}
%%%%%%%%%%%%%%%%%%%%%%%%%%%%%%%%%%%%%%%%%%%%%%%%%%%%%%%%%%%%

Sets are unordered, so an ordered pair must be encoded indirectly.

A standard definition is the Kuratowski ordered pair:

\[
\boxed{
(a,b)
=
\{
\{a\},
\{a,b\}
\}.
}
\]

This representation satisfies the key property

\[
\boxed{
(a,b)
=
(c,d)
\iff
a=c
\land
b=d.
}
\]

%%%%%%%%%%%%%%%%%%%%%%%%%%%%%%%%%%%%%%%%%%%%%%%%%%%%%%%%%%%%
\subsection{Worked Example: Ordered Pair}
\label{subsec:worked-ordered-pair}
%%%%%%%%%%%%%%%%%%%%%%%%%%%%%%%%%%%%%%%%%%%%%%%%%%%%%%%%%%%%

\begin{workedexamplebox}{Encoding \((a,b)\)}

By definition,

\[
(a,b)
=
\{
\{a\},
\{a,b\}
\}.
\]

Thus the ordered information is encoded using nesting.

If

\[
a=b,
\]

then

\[
(a,a)
=
\{
\{a\},
\{a\}
\}
=
\{\{a\}\}.
\]

\begin{answerbox}
\[
\boxed{
(a,b)
=
\{
\{a\},
\{a,b\}
\}.
}
\]
\end{answerbox}

\end{workedexamplebox}

%%%%%%%%%%%%%%%%%%%%%%%%%%%%%%%%%%%%%%%%%%%%%%%%%%%%%%%%%%%%
\subsection{Cartesian Products}
\label{subsec:axiomatic-cartesian-products}
%%%%%%%%%%%%%%%%%%%%%%%%%%%%%%%%%%%%%%%%%%%%%%%%%%%%%%%%%%%%

The Cartesian product is

\[
\boxed{
A\times B
=
\{
(a,b):
a\in A,
b\in B
\}.
}
\]

Because ordered pairs are sets, the Cartesian product can itself be
constructed as a set using the ZF axioms.

One may bound its elements inside a suitable power-set construction and
then apply Separation.

%%%%%%%%%%%%%%%%%%%%%%%%%%%%%%%%%%%%%%%%%%%%%%%%%%%%%%%%%%%%
\subsection{Relations as Sets}
\label{subsec:relations-as-sets}
%%%%%%%%%%%%%%%%%%%%%%%%%%%%%%%%%%%%%%%%%%%%%%%%%%%%%%%%%%%%

A binary relation from \(A\) to \(B\) is a subset

\[
\boxed{
R
\subseteq
A\times B.
}
\]

Thus relations are sets of ordered pairs.

For example,

\[
aRb
\]

means

\[
\boxed{
(a,b)\in R.
}
\]

%%%%%%%%%%%%%%%%%%%%%%%%%%%%%%%%%%%%%%%%%%%%%%%%%%%%%%%%%%%%
\subsection{Functions as Sets}
\label{subsec:functions-as-sets}
%%%%%%%%%%%%%%%%%%%%%%%%%%%%%%%%%%%%%%%%%%%%%%%%%%%%%%%%%%%%

A function

\[
f:A\to B
\]

can be represented as a relation

\[
f
\subseteq
A\times B
\]

such that

\[
\boxed{
\forall x\in A
\exists!y\in B
\,
((x,y)\in f).
}
\]

Thus a function is itself a set.

%%%%%%%%%%%%%%%%%%%%%%%%%%%%%%%%%%%%%%%%%%%%%%%%%%%%%%%%%%%%
\subsection{Set-Theoretic Reduction of Mathematical Objects}
\label{subsec:set-theoretic-reduction}
%%%%%%%%%%%%%%%%%%%%%%%%%%%%%%%%%%%%%%%%%%%%%%%%%%%%%%%%%%%%

Within ZFC, many mathematical objects can be represented entirely using
sets.

For example:

\[
\boxed{
\text{ordered pairs}
\rightarrow
\text{sets},
}
\]

\[
\boxed{
\text{relations}
\rightarrow
\text{sets of ordered pairs},
}
\]

\[
\boxed{
\text{functions}
\rightarrow
\text{relations},
}
\]

and

\[
\boxed{
\text{numbers}
\rightarrow
\text{special sets}.
}
\]

This is one reason set theory serves as a foundational language.

%%%%%%%%%%%%%%%%%%%%%%%%%%%%%%%%%%%%%%%%%%%%%%%%%%%%%%%%%%%%
\subsection{Transitive Sets}
\label{subsec:transitive-sets}
%%%%%%%%%%%%%%%%%%%%%%%%%%%%%%%%%%%%%%%%%%%%%%%%%%%%%%%%%%%%

\begin{definitionbox}{Transitive Set}
A set \(A\) is transitive if every element of an element of \(A\) is also
an element of \(A\).

Equivalently,

\[
\boxed{
x\in y\in A
\Longrightarrow
x\in A.
}
\]

Another equivalent form is

\[
\boxed{
\bigcup A
\subseteq
A.
}
\]
\end{definitionbox}

Von Neumann ordinals are transitive sets.

%%%%%%%%%%%%%%%%%%%%%%%%%%%%%%%%%%%%%%%%%%%%%%%%%%%%%%%%%%%%
\subsection{Ordinals}
\label{subsec:ordinals}
%%%%%%%%%%%%%%%%%%%%%%%%%%%%%%%%%%%%%%%%%%%%%%%%%%%%%%%%%%%%

A von Neumann ordinal is a transitive set whose elements are also
transitive and are well ordered by membership.

Conceptually,

\[
\boxed{
\alpha
=
\{\beta:\beta<\alpha\}.
}
\]

Thus each ordinal is the set of all smaller ordinals.

For finite ordinals this reproduces the natural numbers.

%%%%%%%%%%%%%%%%%%%%%%%%%%%%%%%%%%%%%%%%%%%%%%%%%%%%%%%%%%%%
\subsection{Successor Ordinals}
\label{subsec:successor-ordinals}
%%%%%%%%%%%%%%%%%%%%%%%%%%%%%%%%%%%%%%%%%%%%%%%%%%%%%%%%%%%%

The successor of ordinal \(\alpha\) is

\[
\boxed{
\alpha+1
=
\alpha\cup\{\alpha\}.
}
\]

Thus

\[
0=\varnothing,
\]

\[
1=0\cup\{0\},
\]

\[
2=1\cup\{1\},
\]

and so on.

The notation matches the natural-number successor construction.

%%%%%%%%%%%%%%%%%%%%%%%%%%%%%%%%%%%%%%%%%%%%%%%%%%%%%%%%%%%%
\subsection{Limit Ordinals}
\label{subsec:limit-ordinals}
%%%%%%%%%%%%%%%%%%%%%%%%%%%%%%%%%%%%%%%%%%%%%%%%%%%%%%%%%%%%

A nonzero ordinal that is not a successor ordinal is called a limit
ordinal.

The first infinite ordinal

\[
\boxed{
\omega
}
\]

is a limit ordinal.

Indeed,

\[
\omega
\]

has no largest finite ordinal.

Conceptually,

\[
\boxed{
\omega
=
\sup
\{0,1,2,\ldots\}.
}
\]

%%%%%%%%%%%%%%%%%%%%%%%%%%%%%%%%%%%%%%%%%%%%%%%%%%%%%%%%%%%%
\subsection{Ordinal Order}
\label{subsec:ordinal-order}
%%%%%%%%%%%%%%%%%%%%%%%%%%%%%%%%%%%%%%%%%%%%%%%%%%%%%%%%%%%%

For ordinals \(\alpha\) and \(\beta\),

\[
\boxed{
\alpha<\beta
\iff
\alpha\in\beta.
}
\]

Thus the membership relation itself gives the ordinal ordering.

Every pair of ordinals is comparable:

\[
\boxed{
\alpha\in\beta,
\quad
\alpha=\beta,
\quad
\text{or}
\quad
\beta\in\alpha.
}
\]

%%%%%%%%%%%%%%%%%%%%%%%%%%%%%%%%%%%%%%%%%%%%%%%%%%%%%%%%%%%%
\subsection{Ordinal Arithmetic}
\label{subsec:ordinal-arithmetic}
%%%%%%%%%%%%%%%%%%%%%%%%%%%%%%%%%%%%%%%%%%%%%%%%%%%%%%%%%%%%

Ordinal addition, multiplication, and exponentiation extend finite
arithmetic into the transfinite.

However, ordinal arithmetic is not generally commutative.

For example,

\[
\boxed{
1+\omega
=
\omega,
}
\]

while

\[
\boxed{
\omega+1
>
\omega.
}
\]

Thus

\[
\boxed{
1+\omega
\neq
\omega+1.
}
\]

%%%%%%%%%%%%%%%%%%%%%%%%%%%%%%%%%%%%%%%%%%%%%%%%%%%%%%%%%%%%
\subsection{Why Ordinal Addition Is Not Commutative}
\label{subsec:ordinal-addition-noncommutative}
%%%%%%%%%%%%%%%%%%%%%%%%%%%%%%%%%%%%%%%%%%%%%%%%%%%%%%%%%%%%

Ordinal arithmetic encodes order type.

Appending one point before an infinite sequence does not change the order
type of

\[
\omega.
\]

Thus

\[
1+\omega=\omega.
\]

Appending one point after the infinite sequence produces a new largest
element, giving

\[
\omega+1.
\]

%%%%%%%%%%%%%%%%%%%%%%%%%%%%%%%%%%%%%%%%%%%%%%%%%%%%%%%%%%%%
\subsection{Transfinite Induction}
\label{subsec:transfinite-induction}
%%%%%%%%%%%%%%%%%%%%%%%%%%%%%%%%%%%%%%%%%%%%%%%%%%%%%%%%%%%%

Transfinite induction generalizes ordinary induction to ordinals.

To prove

\[
\boxed{
P(\alpha)
}
\]

for every ordinal \(\alpha\), it is sufficient to show that

\[
\boxed{
\left[
\forall\beta<\alpha
\,
P(\beta)
\right]
\Longrightarrow
P(\alpha).
}
\]

Then \(P(\alpha)\) holds for all ordinals.

%%%%%%%%%%%%%%%%%%%%%%%%%%%%%%%%%%%%%%%%%%%%%%%%%%%%%%%%%%%%
\subsection{Three-Stage Form of Transfinite Induction}
\label{subsec:three-stage-transfinite-induction}
%%%%%%%%%%%%%%%%%%%%%%%%%%%%%%%%%%%%%%%%%%%%%%%%%%%%%%%%%%%%

Transfinite induction is often presented as:

\begin{enumerate}
    \item prove the base case;
    \item prove the successor step;
    \item prove the limit-ordinal step.
\end{enumerate}

Schematically:

\[
\boxed{
P(0),
}
\]

\[
\boxed{
P(\alpha)
\Longrightarrow
P(\alpha+1),
}
\]

and for limit \(\lambda\),

\[
\boxed{
\left[
\forall\beta<\lambda
\,
P(\beta)
\right]
\Longrightarrow
P(\lambda).
}
\]

%%%%%%%%%%%%%%%%%%%%%%%%%%%%%%%%%%%%%%%%%%%%%%%%%%%%%%%%%%%%
\subsection{Transfinite Recursion}
\label{subsec:transfinite-recursion}
%%%%%%%%%%%%%%%%%%%%%%%%%%%%%%%%%%%%%%%%%%%%%%%%%%%%%%%%%%%%

Transfinite recursion defines objects stage by stage along the ordinals.

Conceptually,

\[
\boxed{
F(\alpha)
=
G
\left(
F\restriction\alpha
\right),
}
\]

where the value at stage \(\alpha\) depends on earlier values

\[
F(\beta),
\qquad
\beta<\alpha.
\]

This principle constructs the cumulative hierarchy and many other
transfinite objects.

%%%%%%%%%%%%%%%%%%%%%%%%%%%%%%%%%%%%%%%%%%%%%%%%%%%%%%%%%%%%
\subsection{Equipotence}
\label{subsec:equipotence}
%%%%%%%%%%%%%%%%%%%%%%%%%%%%%%%%%%%%%%%%%%%%%%%%%%%%%%%%%%%%

Two sets \(A\) and \(B\) have the same cardinality if there exists a
bijection

\[
f:A\to B.
\]

Write

\[
\boxed{
A\sim B
}
\]

or

\[
\boxed{
|A|=|B|.
}
\]

This relation is called equipotence.

%%%%%%%%%%%%%%%%%%%%%%%%%%%%%%%%%%%%%%%%%%%%%%%%%%%%%%%%%%%%
\subsection{Cardinal Numbers}
\label{subsec:cardinal-numbers-set-theory}
%%%%%%%%%%%%%%%%%%%%%%%%%%%%%%%%%%%%%%%%%%%%%%%%%%%%%%%%%%%%

Cardinal numbers measure size up to bijection.

A cardinal may be represented by a specially chosen ordinal that is not in
bijection with any smaller ordinal.

Thus cardinals can be realized as initial ordinals.

Finite cardinals coincide with the ordinary natural numbers.

%%%%%%%%%%%%%%%%%%%%%%%%%%%%%%%%%%%%%%%%%%%%%%%%%%%%%%%%%%%%
\subsection{Finite Cardinality}
\label{subsec:finite-cardinality-set-theory}
%%%%%%%%%%%%%%%%%%%%%%%%%%%%%%%%%%%%%%%%%%%%%%%%%%%%%%%%%%%%

A set \(A\) is finite if there exists

\[
n\in\omega
\]

such that

\[
\boxed{
A\sim n.
}
\]

Then

\[
|A|=n.
\]

Otherwise \(A\) is infinite.

%%%%%%%%%%%%%%%%%%%%%%%%%%%%%%%%%%%%%%%%%%%%%%%%%%%%%%%%%%%%
\subsection{Countably Infinite Sets}
\label{subsec:countably-infinite-axiomatic}
%%%%%%%%%%%%%%%%%%%%%%%%%%%%%%%%%%%%%%%%%%%%%%%%%%%%%%%%%%%%

A set \(A\) is countably infinite if

\[
\boxed{
A\sim\omega.
}
\]

Equivalently, its elements can be listed as

\[
\boxed{
a_0,a_1,a_2,\ldots.
}
\]

Examples include

\[
\boxed{
\mathbb{Z}
}
\]

and

\[
\boxed{
\mathbb{Q}.
}
\]

%%%%%%%%%%%%%%%%%%%%%%%%%%%%%%%%%%%%%%%%%%%%%%%%%%%%%%%%%%%%
\subsection{Cantor's Theorem}
\label{subsec:cantors-theorem-set-theory}
%%%%%%%%%%%%%%%%%%%%%%%%%%%%%%%%%%%%%%%%%%%%%%%%%%%%%%%%%%%%

\begin{theorem}
For every set \(A\),

\[
\boxed{
|A|
<
|\mathcal{P}(A)|.
}
\]
\end{theorem}

There is always an injection

\[
A
\to
\mathcal{P}(A)
\]

given by

\[
a\mapsto\{a\}.
\]

Cantor's diagonal argument shows that no surjection

\[
A\to\mathcal{P}(A)
\]

exists.

%%%%%%%%%%%%%%%%%%%%%%%%%%%%%%%%%%%%%%%%%%%%%%%%%%%%%%%%%%%%
\subsection{Proof of Cantor's Theorem}
\label{subsec:proof-cantor-theorem}
%%%%%%%%%%%%%%%%%%%%%%%%%%%%%%%%%%%%%%%%%%%%%%%%%%%%%%%%%%%%

Suppose

\[
f:A\to\mathcal{P}(A)
\]

were surjective.

Define

\[
\boxed{
D
=
\{x\in A:x\notin f(x)\}.
}
\]

Because

\[
D\subseteq A,
\]

we have

\[
D\in\mathcal{P}(A).
\]

Surjectivity gives some

\[
d\in A
\]

such that

\[
f(d)=D.
\]

Then

\[
d\in D
\iff
d\notin f(d)
\iff
d\notin D.
\]

Contradiction.

Therefore no such surjection exists.

%%%%%%%%%%%%%%%%%%%%%%%%%%%%%%%%%%%%%%%%%%%%%%%%%%%%%%%%%%%%
\subsection{Worked Example: Power Sets Are Larger}
\label{subsec:worked-cantor-power-set}
%%%%%%%%%%%%%%%%%%%%%%%%%%%%%%%%%%%%%%%%%%%%%%%%%%%%%%%%%%%%

\begin{workedexamplebox}{Finite Illustration of Cantor's Theorem}

Let

\[
A=\{1,2,3\}.
\]

Then

\[
|A|=3.
\]

Its power set has

\[
2^3=8
\]

elements.

Thus

\[
\begin{answerbox}
\[
\boxed{
|\mathcal{P}(A)|
=
8
>
3
=
|A|.
}
\]
\end{answerbox}

Cantor's theorem says the same strict inequality holds even for infinite
sets.

\end{workedexamplebox}

%%%%%%%%%%%%%%%%%%%%%%%%%%%%%%%%%%%%%%%%%%%%%%%%%%%%%%%%%%%%
\subsection{No Universal Set}
\label{subsec:no-universal-set}
%%%%%%%%%%%%%%%%%%%%%%%%%%%%%%%%%%%%%%%%%%%%%%%%%%%%%%%%%%%%

Suppose there were a set \(V\) containing every set.

Then

\[
\mathcal{P}(V)
\]

would also be a set.

Because every subset of \(V\) is itself a set, we would have

\[
\mathcal{P}(V)
\subseteq
V.
\]

This would imply

\[
|\mathcal{P}(V)|
\leq
|V|,
\]

contradicting Cantor's theorem.

Thus in ZF there is no set of all sets.

%%%%%%%%%%%%%%%%%%%%%%%%%%%%%%%%%%%%%%%%%%%%%%%%%%%%%%%%%%%%
\subsection{The Continuum}
\label{subsec:the-continuum}
%%%%%%%%%%%%%%%%%%%%%%%%%%%%%%%%%%%%%%%%%%%%%%%%%%%%%%%%%%%%

The cardinality of the real numbers is called the cardinality of the
continuum.

One writes

\[
\boxed{
\mathfrak{c}
=
|\mathbb{R}|.
}
\]

There is a bijective correspondence between subsets of \(\omega\) and
binary sequences, and these are closely related to real-number encodings.

Thus

\[
\boxed{
\mathfrak{c}
=
2^{\aleph_0}
=
|\mathcal{P}(\omega)|.
}
\]

%%%%%%%%%%%%%%%%%%%%%%%%%%%%%%%%%%%%%%%%%%%%%%%%%%%%%%%%%%%%
\subsection{Aleph Numbers}
\label{subsec:aleph-numbers}
%%%%%%%%%%%%%%%%%%%%%%%%%%%%%%%%%%%%%%%%%%%%%%%%%%%%%%%%%%%%

Infinite cardinal numbers are commonly indexed by ordinals:

\[
\boxed{
\aleph_0,
\aleph_1,
\aleph_2,
\ldots.
}
\]

Here

\[
\aleph_0
\]

is the cardinality of \(\omega\).

The next larger well-orderable cardinal is

\[
\aleph_1.
\]

The sequence continues transfinitely.

%%%%%%%%%%%%%%%%%%%%%%%%%%%%%%%%%%%%%%%%%%%%%%%%%%%%%%%%%%%%
\subsection{Continuum Hypothesis}
\label{subsec:continuum-hypothesis}
%%%%%%%%%%%%%%%%%%%%%%%%%%%%%%%%%%%%%%%%%%%%%%%%%%%%%%%%%%%%

The Continuum Hypothesis, abbreviated CH, asserts that there is no
cardinality strictly between

\[
\aleph_0
\]

and

\[
2^{\aleph_0}.
\]

Equivalently,

\[
\boxed{
2^{\aleph_0}
=
\aleph_1.
}
\]

CH is not provable or disprovable from ZFC, assuming ZFC is consistent.

%%%%%%%%%%%%%%%%%%%%%%%%%%%%%%%%%%%%%%%%%%%%%%%%%%%%%%%%%%%%
\subsection{Generalized Continuum Hypothesis}
\label{subsec:generalized-continuum-hypothesis}
%%%%%%%%%%%%%%%%%%%%%%%%%%%%%%%%%%%%%%%%%%%%%%%%%%%%%%%%%%%%

The Generalized Continuum Hypothesis, abbreviated GCH, asserts that for
every infinite cardinal \(\kappa\),

\[
\boxed{
2^\kappa
=
\kappa^+,
}
\]

where

\[
\kappa^+
\]

is the next larger cardinal.

GCH is stronger than CH.

%%%%%%%%%%%%%%%%%%%%%%%%%%%%%%%%%%%%%%%%%%%%%%%%%%%%%%%%%%%%
\subsection{Cardinal Arithmetic}
\label{subsec:cardinal-arithmetic}
%%%%%%%%%%%%%%%%%%%%%%%%%%%%%%%%%%%%%%%%%%%%%%%%%%%%%%%%%%%%

For cardinalities

\[
\kappa=|A|,
\qquad
\lambda=|B|,
\]

cardinal addition corresponds to disjoint unions:

\[
\boxed{
\kappa+\lambda
=
|A\sqcup B|.
}
\]

Cardinal multiplication corresponds to Cartesian products:

\[
\boxed{
\kappa\lambda
=
|A\times B|.
}
\]

Cardinal exponentiation corresponds to function sets:

\[
\boxed{
\lambda^\kappa
=
|B^A|,
}
\]

where

\[
B^A
=
\{f:A\to B\}.
\]

%%%%%%%%%%%%%%%%%%%%%%%%%%%%%%%%%%%%%%%%%%%%%%%%%%%%%%%%%%%%
\subsection{Infinite Cardinal Arithmetic}
\label{subsec:infinite-cardinal-arithmetic}
%%%%%%%%%%%%%%%%%%%%%%%%%%%%%%%%%%%%%%%%%%%%%%%%%%%%%%%%%%%%

Assuming Choice, for infinite cardinals \(\kappa,\lambda\) with at least one
nonzero,

\[
\boxed{
\kappa+\lambda
=
\kappa\lambda
=
\max(\kappa,\lambda)
}
\]

under the standard infinite-cardinal hypotheses.

For example,

\[
\boxed{
\aleph_0+\aleph_0
=
\aleph_0,
}
\]

and

\[
\boxed{
\aleph_0\cdot\aleph_0
=
\aleph_0.
}
\]

%%%%%%%%%%%%%%%%%%%%%%%%%%%%%%%%%%%%%%%%%%%%%%%%%%%%%%%%%%%%
\subsection{Well-Ordered Sets}
\label{subsec:well-ordered-sets}
%%%%%%%%%%%%%%%%%%%%%%%%%%%%%%%%%%%%%%%%%%%%%%%%%%%%%%%%%%%%

A linearly ordered set \((A,<)\) is well ordered if every nonempty subset

\[
B\subseteq A
\]

has a least element.

Thus

\[
\boxed{
\forall B
\subseteq
A,
\quad
B\neq\varnothing
\Longrightarrow
B
\text{ has a least element}.
}
\]

The natural numbers with their usual order are well ordered.

%%%%%%%%%%%%%%%%%%%%%%%%%%%%%%%%%%%%%%%%%%%%%%%%%%%%%%%%%%%%
\subsection{Well-Ordering Theorem}
\label{subsec:well-ordering-theorem}
%%%%%%%%%%%%%%%%%%%%%%%%%%%%%%%%%%%%%%%%%%%%%%%%%%%%%%%%%%%%

The Well-Ordering Theorem states:

\[
\boxed{
\text{Every set can be well ordered.}
}
\]

Within ZF, the Well-Ordering Theorem is equivalent to the Axiom of Choice.

Thus Choice can be reformulated as an ordering principle.

%%%%%%%%%%%%%%%%%%%%%%%%%%%%%%%%%%%%%%%%%%%%%%%%%%%%%%%%%%%%
\subsection{Zorn's Lemma}
\label{subsec:zorns-lemma}
%%%%%%%%%%%%%%%%%%%%%%%%%%%%%%%%%%%%%%%%%%%%%%%%%%%%%%%%%%%%

\begin{theorem}
Let \(P\) be a partially ordered set.

If every chain in \(P\) has an upper bound in \(P\), then \(P\) has a
maximal element.
\end{theorem}

Zorn's Lemma is equivalent to the Axiom of Choice over ZF.

It is one of the most frequently used forms of Choice in algebra and
analysis.

%%%%%%%%%%%%%%%%%%%%%%%%%%%%%%%%%%%%%%%%%%%%%%%%%%%%%%%%%%%%
\subsection{Maximal Versus Maximum}
\label{subsec:maximal-versus-maximum}
%%%%%%%%%%%%%%%%%%%%%%%%%%%%%%%%%%%%%%%%%%%%%%%%%%%%%%%%%%%%

A maximum element \(m\) satisfies

\[
\boxed{
x\leq m
}
\]

for every \(x\) in the poset.

A maximal element \(m\) satisfies that there is no strictly larger
comparable element:

\[
\boxed{
m<x
}
\]

for no \(x\) in the poset.

A partially ordered set may have several maximal elements but no maximum.

%%%%%%%%%%%%%%%%%%%%%%%%%%%%%%%%%%%%%%%%%%%%%%%%%%%%%%%%%%%%
\subsection{Worked Example: Maximal Elements}
\label{subsec:worked-maximal-elements}
%%%%%%%%%%%%%%%%%%%%%%%%%%%%%%%%%%%%%%%%%%%%%%%%%%%%%%%%%%%%

\begin{workedexamplebox}{Maximal Does Not Mean Maximum}

Consider subsets

\[
A=\{1\},
\qquad
B=\{2\},
\]

ordered by inclusion.

Neither contains the other.

Thus both are maximal within

\[
\{A,B\},
\]

but there is no maximum element.

\begin{answerbox}
\[
\boxed{
A
\text{ and }
B
\text{ are maximal, but neither is maximum.}
}
\]
\end{answerbox}

\end{workedexamplebox}

%%%%%%%%%%%%%%%%%%%%%%%%%%%%%%%%%%%%%%%%%%%%%%%%%%%%%%%%%%%%
\subsection{Typical Use of Zorn's Lemma}
\label{subsec:typical-zorn-use}
%%%%%%%%%%%%%%%%%%%%%%%%%%%%%%%%%%%%%%%%%%%%%%%%%%%%%%%%%%%%

A standard Zorn argument follows this pattern:

\begin{enumerate}
    \item define a partially ordered collection of partial objects;
    \item order them by extension or inclusion;
    \item show every chain has an upper bound;
    \item apply Zorn's Lemma;
    \item obtain a maximal object;
    \item prove maximality implies the desired property.
\end{enumerate}

This method is used in many existence proofs.

%%%%%%%%%%%%%%%%%%%%%%%%%%%%%%%%%%%%%%%%%%%%%%%%%%%%%%%%%%%%
\subsection{Example: Bases of Vector Spaces}
\label{subsec:zorn-vector-space-basis}
%%%%%%%%%%%%%%%%%%%%%%%%%%%%%%%%%%%%%%%%%%%%%%%%%%%%%%%%%%%%

Let \(V\) be a vector space.

Consider all linearly independent subsets of \(V\), ordered by inclusion.

The union of a chain of linearly independent sets is linearly independent.

Therefore every chain has an upper bound.

Zorn's Lemma yields a maximal linearly independent set.

Such a set is a basis.

Thus Choice implies every vector space has a basis.

%%%%%%%%%%%%%%%%%%%%%%%%%%%%%%%%%%%%%%%%%%%%%%%%%%%%%%%%%%%%
\subsection{Choice Principles}
\label{subsec:choice-principles}
%%%%%%%%%%%%%%%%%%%%%%%%%%%%%%%%%%%%%%%%%%%%%%%%%%%%%%%%%%%%

Over ZF, the following are equivalent:

\[
\boxed{
\text{Axiom of Choice},
}
\]

\[
\boxed{
\text{Well-Ordering Theorem},
}
\]

and

\[
\boxed{
\text{Zorn's Lemma}.
}
\]

Many other equivalent formulations exist.

Different areas of mathematics often use whichever form is most natural.

%%%%%%%%%%%%%%%%%%%%%%%%%%%%%%%%%%%%%%%%%%%%%%%%%%%%%%%%%%%%
\subsection{Filters}
\label{subsec:set-theoretic-filters}
%%%%%%%%%%%%%%%%%%%%%%%%%%%%%%%%%%%%%%%%%%%%%%%%%%%%%%%%%%%%

Let \(X\) be a set.

A filter \(\mathcal{F}\) on \(X\) is a family of subsets satisfying:

\[
\boxed{
X\in\mathcal{F},
}
\]

\[
\boxed{
\varnothing\notin\mathcal{F},
}
\]

\[
A,B\in\mathcal{F}
\Longrightarrow
A\cap B\in\mathcal{F},
\]

and

\[
A\in\mathcal{F},
\quad
A\subseteq B\subseteq X
\Longrightarrow
B\in\mathcal{F}.
\]

Filters formalize collections of ``large'' subsets.

%%%%%%%%%%%%%%%%%%%%%%%%%%%%%%%%%%%%%%%%%%%%%%%%%%%%%%%%%%%%
\subsection{Ultrafilters}
\label{subsec:ultrafilters-preview}
%%%%%%%%%%%%%%%%%%%%%%%%%%%%%%%%%%%%%%%%%%%%%%%%%%%%%%%%%%%%

An ultrafilter is a maximal proper filter.

Equivalently, for every subset \(A\subseteq X\), an ultrafilter
\(\mathcal{U}\) satisfies exactly one of

\[
\boxed{
A\in\mathcal{U}
}
\]

or

\[
\boxed{
X\setminus A\in\mathcal{U}.
}
\]

Ultrafilters play major roles in topology and model theory.

Their existence in general is connected with choice principles.

%%%%%%%%%%%%%%%%%%%%%%%%%%%%%%%%%%%%%%%%%%%%%%%%%%%%%%%%%%%%
\subsection{Transitive Closure}
\label{subsec:transitive-closure-set-theory}
%%%%%%%%%%%%%%%%%%%%%%%%%%%%%%%%%%%%%%%%%%%%%%%%%%%%%%%%%%%%

The transitive closure of a set \(A\) is the smallest transitive set
containing \(A\).

Conceptually,

\[
\boxed{
\operatorname{TC}(A)
=
A
\cup
\bigcup A
\cup
\bigcup\bigcup A
\cup
\cdots.
}
\]

It collects elements, elements of elements, and so forth.

%%%%%%%%%%%%%%%%%%%%%%%%%%%%%%%%%%%%%%%%%%%%%%%%%%%%%%%%%%%%
\subsection{Rank}
\label{subsec:set-theoretic-rank}
%%%%%%%%%%%%%%%%%%%%%%%%%%%%%%%%%%%%%%%%%%%%%%%%%%%%%%%%%%%%

The rank of a set measures how high it appears in the cumulative hierarchy.

It is defined recursively by

\[
\boxed{
\operatorname{rank}(x)
=
\sup
\{
\operatorname{rank}(y)+1:
y\in x
\}.
}
\]

The empty set has rank

\[
\boxed{
0.
}
\]

%%%%%%%%%%%%%%%%%%%%%%%%%%%%%%%%%%%%%%%%%%%%%%%%%%%%%%%%%%%%
\subsection{Cumulative Hierarchy}
\label{subsec:cumulative-hierarchy}
%%%%%%%%%%%%%%%%%%%%%%%%%%%%%%%%%%%%%%%%%%%%%%%%%%%%%%%%%%%%

Define stages \(V_\alpha\) recursively.

Start with

\[
\boxed{
V_0
=
\varnothing.
}
\]

For successor ordinals,

\[
\boxed{
V_{\alpha+1}
=
\mathcal{P}(V_\alpha).
}
\]

For limit ordinals \(\lambda\),

\[
\boxed{
V_\lambda
=
\bigcup_{\alpha<\lambda}
V_\alpha.
}
\]

%%%%%%%%%%%%%%%%%%%%%%%%%%%%%%%%%%%%%%%%%%%%%%%%%%%%%%%%%%%%
\subsection{First Levels of the Hierarchy}
\label{subsec:first-levels-cumulative-hierarchy}
%%%%%%%%%%%%%%%%%%%%%%%%%%%%%%%%%%%%%%%%%%%%%%%%%%%%%%%%%%%%

We have

\[
V_0
=
\varnothing.
\]

Then

\[
V_1
=
\mathcal{P}(\varnothing)
=
\{\varnothing\}.
\]

Next,

\[
V_2
=
\mathcal{P}(V_1)
=
\{
\varnothing,
\{\varnothing\}
\}.
\]

Then

\[
V_3
=
\mathcal{P}(V_2).
\]

Each stage contains sets constructed from earlier stages.

%%%%%%%%%%%%%%%%%%%%%%%%%%%%%%%%%%%%%%%%%%%%%%%%%%%%%%%%%%%%
\subsection{The Universe \(V\)}
\label{subsec:set-theoretic-universe-v}
%%%%%%%%%%%%%%%%%%%%%%%%%%%%%%%%%%%%%%%%%%%%%%%%%%%%%%%%%%%%

The cumulative set-theoretic universe is described informally by

\[
\boxed{
V
=
\bigcup_{\alpha\in\operatorname{Ord}}
V_\alpha.
}
\]

This is not a set inside ZFC.

It is a proper class representing all sets in the cumulative hierarchy.

Every set appears at some stage \(V_\alpha\).

%%%%%%%%%%%%%%%%%%%%%%%%%%%%%%%%%%%%%%%%%%%%%%%%%%%%%%%%%%%%
\subsection{Why the Hierarchy Avoids Self-Membership}
\label{subsec:hierarchy-no-self-membership}
%%%%%%%%%%%%%%%%%%%%%%%%%%%%%%%%%%%%%%%%%%%%%%%%%%%%%%%%%%%%

If

\[
x\in V_{\alpha+1},
\]

then

\[
x\subseteq V_\alpha.
\]

The elements of \(x\) therefore lie at earlier stages.

Thus membership moves downward through rank.

This makes structures such as

\[
x\in x
\]

incompatible with the hierarchy.

%%%%%%%%%%%%%%%%%%%%%%%%%%%%%%%%%%%%%%%%%%%%%%%%%%%%%%%%%%%%
\subsection{Proper Classes}
\label{subsec:proper-classes}
%%%%%%%%%%%%%%%%%%%%%%%%%%%%%%%%%%%%%%%%%%%%%%%%%%%%%%%%%%%%

Some mathematically meaningful collections are too large to be sets.

Such collections are called proper classes.

Examples include:

\[
\boxed{
\text{the class of all sets},
}
\]

and

\[
\boxed{
\text{the class of all ordinals}.
}
\]

In ZFC, proper classes are usually treated metamathematically rather than as
objects in the first-order domain.

%%%%%%%%%%%%%%%%%%%%%%%%%%%%%%%%%%%%%%%%%%%%%%%%%%%%%%%%%%%%
\subsection{The Class of All Ordinals}
\label{subsec:class-all-ordinals}
%%%%%%%%%%%%%%%%%%%%%%%%%%%%%%%%%%%%%%%%%%%%%%%%%%%%%%%%%%%%

Let

\[
\operatorname{Ord}
\]

denote the collection of all ordinals.

It cannot be a set.

If it were a set of all ordinals, it would itself be an ordinal larger than
every ordinal it contains, leading to a contradiction.

This is related to the Burali--Forti paradox.

%%%%%%%%%%%%%%%%%%%%%%%%%%%%%%%%%%%%%%%%%%%%%%%%%%%%%%%%%%%%
\subsection{Burali--Forti Paradox}
\label{subsec:burali-forti-paradox}
%%%%%%%%%%%%%%%%%%%%%%%%%%%%%%%%%%%%%%%%%%%%%%%%%%%%%%%%%%%%

Suppose the collection of all ordinals formed a set

\[
\Omega.
\]

Because ordinals are well ordered by membership, \(\Omega\) would itself
behave as an ordinal.

Then

\[
\Omega
\]

would have to be greater than every ordinal, including itself.

This is impossible.

Therefore the ordinals form a proper class rather than a set.

%%%%%%%%%%%%%%%%%%%%%%%%%%%%%%%%%%%%%%%%%%%%%%%%%%%%%%%%%%%%
\subsection{Class of All Sets}
\label{subsec:class-of-all-sets}
%%%%%%%%%%%%%%%%%%%%%%%%%%%%%%%%%%%%%%%%%%%%%%%%%%%%%%%%%%%%

The collection of all sets is commonly denoted informally by

\[
\boxed{
V.
}
\]

It is a proper class.

If \(V\) were a set, Cantor's theorem would imply

\[
|\mathcal{P}(V)|
>
|V|.
\]

But every subset of \(V\) would itself be a set and hence belong to \(V\),
creating a contradiction.

%%%%%%%%%%%%%%%%%%%%%%%%%%%%%%%%%%%%%%%%%%%%%%%%%%%%%%%%%%%%
\subsection{Set Theory as a First-Order Theory}
\label{subsec:set-theory-first-order-theory}
%%%%%%%%%%%%%%%%%%%%%%%%%%%%%%%%%%%%%%%%%%%%%%%%%%%%%%%%%%%%

ZFC is usually formulated as a first-order theory.

Its variables range over sets.

The only primitive nonlogical relation is

\[
\boxed{
\in.
}
\]

Thus statements about ordered pairs, numbers, functions, and mathematical
spaces are ultimately encoded using sets and membership.

%%%%%%%%%%%%%%%%%%%%%%%%%%%%%%%%%%%%%%%%%%%%%%%%%%%%%%%%%%%%
\subsection{Models of Set Theory}
\label{subsec:models-of-set-theory}
%%%%%%%%%%%%%%%%%%%%%%%%%%%%%%%%%%%%%%%%%%%%%%%%%%%%%%%%%%%%

A model of ZF or ZFC is a structure

\[
\mathcal{M}
\]

whose membership relation satisfies the axioms.

Studying such models is central to advanced set theory and model theory.

A statement may hold in one model of ZFC and fail in another if the
statement is independent of ZFC.

%%%%%%%%%%%%%%%%%%%%%%%%%%%%%%%%%%%%%%%%%%%%%%%%%%%%%%%%%%%%
\subsection{Internal and External Perspectives}
\label{subsec:internal-external-set-theory}
%%%%%%%%%%%%%%%%%%%%%%%%%%%%%%%%%%%%%%%%%%%%%%%%%%%%%%%%%%%%

When studying a model \(\mathcal{M}\), one distinguishes:

\[
\boxed{
\text{what \(\mathcal{M}\) believes}
}
\]

from

\[
\boxed{
\text{what is true about \(\mathcal{M}\) externally}.
}
\]

For example, a model may internally regard a set as uncountable even when,
from an external metamathematical perspective, the model itself is
countable.

This distinction appears prominently in model theory.

%%%%%%%%%%%%%%%%%%%%%%%%%%%%%%%%%%%%%%%%%%%%%%%%%%%%%%%%%%%%
\subsection{Relative Consistency}
\label{subsec:relative-consistency}
%%%%%%%%%%%%%%%%%%%%%%%%%%%%%%%%%%%%%%%%%%%%%%%%%%%%%%%%%%%%

Because sufficiently strong formal systems cannot generally prove their
own consistency under standard assumptions, set-theoretic consistency
results are often relative.

A statement such as

\[
\boxed{
\operatorname{Con}(\mathrm{ZFC})
\Longrightarrow
\operatorname{Con}(\mathrm{ZFC}+\varphi)
}
\]

means that if ZFC is consistent, then adjoining \(\varphi\) also yields a
consistent theory.

This does not prove the absolute consistency of ZFC.

%%%%%%%%%%%%%%%%%%%%%%%%%%%%%%%%%%%%%%%%%%%%%%%%%%%%%%%%%%%%
\subsection{Independence}
\label{subsec:set-theoretic-independence}
%%%%%%%%%%%%%%%%%%%%%%%%%%%%%%%%%%%%%%%%%%%%%%%%%%%%%%%%%%%%

A sentence \(\varphi\) is independent of theory \(T\) when neither

\[
\boxed{
T\vdash\varphi
}
\]

nor

\[
\boxed{
T\vdash\neg\varphi
}
\]

holds, assuming \(T\) is consistent in the relevant metatheoretic setting.

Then both

\[
T+\varphi
\]

and

\[
T+\neg\varphi
\]

may be relatively consistent.

%%%%%%%%%%%%%%%%%%%%%%%%%%%%%%%%%%%%%%%%%%%%%%%%%%%%%%%%%%%%
\subsection{Continuum Hypothesis Independence}
\label{subsec:ch-independence}
%%%%%%%%%%%%%%%%%%%%%%%%%%%%%%%%%%%%%%%%%%%%%%%%%%%%%%%%%%%%

The Continuum Hypothesis provides the canonical example.

Gödel showed that, assuming ZF is consistent, adding suitable forms of
Choice and CH does not create inconsistency through the constructible
universe.

Cohen later showed, using forcing, that CH cannot be proved from ZFC if ZFC
is consistent.

Thus CH is independent of ZFC relative to the usual consistency
assumption.

%%%%%%%%%%%%%%%%%%%%%%%%%%%%%%%%%%%%%%%%%%%%%%%%%%%%%%%%%%%%
\subsection{Constructible Universe Preview}
\label{subsec:constructible-universe-preview}
%%%%%%%%%%%%%%%%%%%%%%%%%%%%%%%%%%%%%%%%%%%%%%%%%%%%%%%%%%%%

The constructible universe, denoted

\[
\boxed{
L,
}
\]

is an inner model built through a hierarchy resembling \(V\), but including
only sets that are definable from earlier stages.

Conceptually,

\[
\boxed{
L_0=\varnothing,
}
\]

and later stages contain subsets definable over earlier stages.

The constructible universe satisfies strong structural principles,
including Choice and GCH.

%%%%%%%%%%%%%%%%%%%%%%%%%%%%%%%%%%%%%%%%%%%%%%%%%%%%%%%%%%%%
\subsection{Forcing Preview}
\label{subsec:forcing-preview}
%%%%%%%%%%%%%%%%%%%%%%%%%%%%%%%%%%%%%%%%%%%%%%%%%%%%%%%%%%%%

Forcing is a method for constructing models of set theory with new sets.

It is used to show relative consistency and independence results.

Conceptually, forcing begins with a model \(M\) and adjoins a suitable
generic object \(G\), producing

\[
\boxed{
M[G].
}
\]

The resulting extension satisfies carefully controlled statements.

%%%%%%%%%%%%%%%%%%%%%%%%%%%%%%%%%%%%%%%%%%%%%%%%%%%%%%%%%%%%
\subsection{Why Forcing Matters}
\label{subsec:why-forcing-matters}
%%%%%%%%%%%%%%%%%%%%%%%%%%%%%%%%%%%%%%%%%%%%%%%%%%%%%%%%%%%%

Forcing demonstrates that many questions cannot be settled by the usual ZFC
axioms alone.

It has been used to study:

\[
\boxed{
\text{the Continuum Hypothesis},
}
\]

\[
\boxed{
\text{cardinal arithmetic},
}
\]

\[
\boxed{
\text{combinatorial principles},
}
\]

and

\[
\boxed{
\text{properties of the real line}.
}
\]

It is one of the central techniques of modern set theory.

%%%%%%%%%%%%%%%%%%%%%%%%%%%%%%%%%%%%%%%%%%%%%%%%%%%%%%%%%%%%
\subsection{Large Cardinals Preview}
\label{subsec:large-cardinals-preview}
%%%%%%%%%%%%%%%%%%%%%%%%%%%%%%%%%%%%%%%%%%%%%%%%%%%%%%%%%%%%

Large cardinal axioms assert the existence of extremely large infinite
cardinals with strong structural properties.

Examples include concepts such as:

\[
\boxed{
\text{inaccessible cardinals},
}
\]

\[
\boxed{
\text{measurable cardinals},
}
\]

and stronger notions.

Large cardinal axioms extend the strength of ordinary ZFC.

%%%%%%%%%%%%%%%%%%%%%%%%%%%%%%%%%%%%%%%%%%%%%%%%%%%%%%%%%%%%
\subsection{Inaccessible Cardinals Preview}
\label{subsec:inaccessible-cardinals-preview}
%%%%%%%%%%%%%%%%%%%%%%%%%%%%%%%%%%%%%%%%%%%%%%%%%%%%%%%%%%%%

Very roughly, an inaccessible cardinal \(\kappa\) is an uncountable
regular strong-limit cardinal.

Such a cardinal is closed under many set-theoretic operations below
\(\kappa\).

The hierarchy level

\[
V_\kappa
\]

then behaves in many ways like a small universe of sets.

%%%%%%%%%%%%%%%%%%%%%%%%%%%%%%%%%%%%%%%%%%%%%%%%%%%%%%%%%%%%
\subsection{Why Large Cardinals Matter}
\label{subsec:why-large-cardinals-matter}
%%%%%%%%%%%%%%%%%%%%%%%%%%%%%%%%%%%%%%%%%%%%%%%%%%%%%%%%%%%%

Large cardinal axioms provide:

\[
\boxed{
\text{strong consistency assumptions},
}
\]

\[
\boxed{
\text{structural principles for the set-theoretic universe},
}
\]

and

\[
\boxed{
\text{benchmarks of logical strength}.
}
\]

They also interact deeply with descriptive set theory, inner-model theory,
and determinacy principles.

%%%%%%%%%%%%%%%%%%%%%%%%%%%%%%%%%%%%%%%%%%%%%%%%%%%%%%%%%%%%
\subsection{Set Theory and Mathematical Foundations}
\label{subsec:set-theory-foundations}
%%%%%%%%%%%%%%%%%%%%%%%%%%%%%%%%%%%%%%%%%%%%%%%%%%%%%%%%%%%%

Set theory provides one way to encode most ordinary mathematical objects in
a single foundational universe.

For example:

\[
\boxed{
\mathbb{N}
\rightarrow
\text{sets},
}
\]

\[
\boxed{
\mathbb{Z}
\rightarrow
\text{equivalence classes of pairs of naturals},
}
\]

\[
\boxed{
\mathbb{Q}
\rightarrow
\text{equivalence classes of integer pairs},
}
\]

and

\[
\boxed{
\mathbb{R}
\rightarrow
\text{Dedekind cuts or Cauchy-sequence constructions}.
}
\]

Structures such as groups and topological spaces can likewise be coded as
sets.

%%%%%%%%%%%%%%%%%%%%%%%%%%%%%%%%%%%%%%%%%%%%%%%%%%%%%%%%%%%%
\subsection{Set-Theoretic Universe Versus Mathematical Practice}
\label{subsec:set-universe-vs-practice}
%%%%%%%%%%%%%%%%%%%%%%%%%%%%%%%%%%%%%%%%%%%%%%%%%%%%%%%%%%%%

In ordinary mathematical work, one rarely expands every object into its
set-theoretic encoding.

A mathematician writes

\[
f:X\to Y
\]

without replacing \(f\), \(X\), and \(Y\) by nested sets.

Set theory provides the foundational justification in the background.

Mathematical practice then uses higher-level structures for clarity.

%%%%%%%%%%%%%%%%%%%%%%%%%%%%%%%%%%%%%%%%%%%%%%%%%%%%%%%%%%%%
\subsection{Set Theory and Category Theory}
\label{subsec:set-theory-category-theory}
%%%%%%%%%%%%%%%%%%%%%%%%%%%%%%%%%%%%%%%%%%%%%%%%%%%%%%%%%%%%

Category theory often works with collections that are too large to be sets.

For example, the collection of all sets or all groups is naturally a
proper class.

Foundational treatments may therefore use:

\[
\boxed{
\text{set/class distinctions},
}
\]

\[
\boxed{
\text{universes},
}
\]

or other foundational frameworks.

Set theory and category theory are complementary rather than mutually
exclusive foundations.

%%%%%%%%%%%%%%%%%%%%%%%%%%%%%%%%%%%%%%%%%%%%%%%%%%%%%%%%%%%%
\subsection{Axiomatic Set Theory Workflow}
\label{subsec:axiomatic-set-theory-workflow}
%%%%%%%%%%%%%%%%%%%%%%%%%%%%%%%%%%%%%%%%%%%%%%%%%%%%%%%%%%%%

When analyzing a set-theoretic construction:

\begin{enumerate}
    \item identify the sets already known to exist;
    \item determine which axiom permits the new construction;
    \item use Pairing for small finite collections;
    \item use Union to flatten one membership level;
    \item use Power Set to collect subsets;
    \item use Separation to select a definable subset;
    \item use Replacement to collect definable images;
    \item use Infinity when an infinite inductive set is required;
    \item use Foundation when well-founded membership matters;
    \item identify whether Choice is required;
    \item distinguish sets from proper classes;
    \item specify the relevant ordinal or cardinal when working
          transfinitely;
    \item use transfinite induction for ordinal-indexed proofs;
    \item use transfinite recursion for ordinal-indexed constructions;
    \item verify that a purported comprehension is restricted to an
          existing set;
    \item distinguish theorem statements from added axioms;
    \item distinguish provability from independence;
    \item state any relative consistency assumptions explicitly.
\end{enumerate}

%%%%%%%%%%%%%%%%%%%%%%%%%%%%%%%%%%%%%%%%%%%%%%%%%%%%%%%%%%%%
\subsection{Common Errors}
\label{subsec:axiomatic-set-theory-common-errors}
%%%%%%%%%%%%%%%%%%%%%%%%%%%%%%%%%%%%%%%%%%%%%%%%%%%%%%%%%%%%

\subsubsection{Assuming Every Collection Is a Set}

Collections such as

\[
\boxed{
\text{all sets}
}
\]

and

\[
\boxed{
\text{all ordinals}
}
\]

cannot be sets in ZFC.

They are proper classes.

%%%%%%%%%%%%%%%%%%%%%%%%%%%%%%%%%%%%%%%%%%%%%%%%%%%%%%%%%%%%
\subsubsection{Using Unrestricted Comprehension}

The expression

\[
\{x:P(x)\}
\]

is not automatically legitimate.

A safe ZF construction has form

\[
\boxed{
\{x\in A:P(x)\},
}
\]

where \(A\) is already known to be a set.

%%%%%%%%%%%%%%%%%%%%%%%%%%%%%%%%%%%%%%%%%%%%%%%%%%%%%%%%%%%%
\subsubsection{Confusing Membership and Subset}

The statement

\[
x\in A
\]

means \(x\) is an element of \(A\).

The statement

\[
x\subseteq A
\]

means every element of \(x\) lies in \(A\).

These are completely different relations.

%%%%%%%%%%%%%%%%%%%%%%%%%%%%%%%%%%%%%%%%%%%%%%%%%%%%%%%%%%%%
\subsubsection{Assuming \(x\in x\) Is Allowed in Standard ZFC}

Foundation rules out self-membership.

Standard well-founded set theory does not contain sets satisfying

\[
x\in x.
\]

%%%%%%%%%%%%%%%%%%%%%%%%%%%%%%%%%%%%%%%%%%%%%%%%%%%%%%%%%%%%
\subsubsection{Assuming the Empty Set Is an Element of Every Set}

The empty set is a subset of every set:

\[
\boxed{
\varnothing\subseteq A.
}
\]

But it need not be an element:

\[
\boxed{
\varnothing\in A
}
\]

may be true or false.

%%%%%%%%%%%%%%%%%%%%%%%%%%%%%%%%%%%%%%%%%%%%%%%%%%%%%%%%%%%%
\subsubsection{Confusing \(\{a\}\) With \(a\)}

A singleton set

\[
\{a\}
\]

contains \(a\) as one element.

It is generally not equal to \(a\).

For von Neumann naturals, this distinction remains important.

%%%%%%%%%%%%%%%%%%%%%%%%%%%%%%%%%%%%%%%%%%%%%%%%%%%%%%%%%%%%
\subsubsection{Confusing \(\in\) With Ordinal \(<\) Outside Ordinals}

For ordinals,

\[
\alpha<\beta
\iff
\alpha\in\beta.
\]

This does not mean that membership is an ordinary ordering on arbitrary
sets.

%%%%%%%%%%%%%%%%%%%%%%%%%%%%%%%%%%%%%%%%%%%%%%%%%%%%%%%%%%%%
\subsubsection{Assuming Ordinal Arithmetic Is Commutative}

In general,

\[
\boxed{
\alpha+\beta
\neq
\beta+\alpha.
}
\]

For example,

\[
1+\omega=\omega,
\]

but

\[
\omega+1>\omega.
\]

%%%%%%%%%%%%%%%%%%%%%%%%%%%%%%%%%%%%%%%%%%%%%%%%%%%%%%%%%%%%
\subsubsection{Confusing Ordinals and Cardinals}

Ordinals describe order type.

Cardinals describe size up to bijection.

The same cardinality may support many different order types.

%%%%%%%%%%%%%%%%%%%%%%%%%%%%%%%%%%%%%%%%%%%%%%%%%%%%%%%%%%%%
\subsubsection{Assuming Infinite Sets Behave Like Finite Sets}

For infinite sets it is possible that

\[
A
\subsetneq
B
\]

while

\[
|A|=|B|.
\]

For example,

\[
\mathbb{N}
\]

and the even natural numbers have the same cardinality.

%%%%%%%%%%%%%%%%%%%%%%%%%%%%%%%%%%%%%%%%%%%%%%%%%%%%%%%%%%%%
\subsubsection{Assuming \(|A|=|\mathcal{P}(A)|\) for Infinite Sets}

Cantor's theorem gives

\[
\boxed{
|A|
<
|\mathcal{P}(A)|
}
\]

for every set \(A\), finite or infinite.

%%%%%%%%%%%%%%%%%%%%%%%%%%%%%%%%%%%%%%%%%%%%%%%%%%%%%%%%%%%%
\subsubsection{Confusing Countable With Finite}

A countably infinite set has the same cardinality as

\[
\mathbb{N}
\]

but is not finite.

%%%%%%%%%%%%%%%%%%%%%%%%%%%%%%%%%%%%%%%%%%%%%%%%%%%%%%%%%%%%
\subsubsection{Assuming Choice Is Provable in ZF}

The Axiom of Choice is independent of ZF relative to the usual consistency
assumptions.

ZFC explicitly adds Choice.

%%%%%%%%%%%%%%%%%%%%%%%%%%%%%%%%%%%%%%%%%%%%%%%%%%%%%%%%%%%%
\subsubsection{Confusing Maximal and Maximum}

A maximal element need not dominate every element.

Zorn's Lemma guarantees maximal elements, not necessarily maximum elements.

%%%%%%%%%%%%%%%%%%%%%%%%%%%%%%%%%%%%%%%%%%%%%%%%%%%%%%%%%%%%
\subsubsection{Using Zorn's Lemma Without Checking Chain Upper Bounds}

The key hypothesis is that every chain has an upper bound inside the
partially ordered set under consideration.

Without this, Zorn's Lemma does not apply.

%%%%%%%%%%%%%%%%%%%%%%%%%%%%%%%%%%%%%%%%%%%%%%%%%%%%%%%%%%%%
\subsubsection{Assuming CH Is True or False in ZFC}

The Continuum Hypothesis is independent of ZFC, assuming ZFC is consistent.

Therefore ZFC alone does not decide CH.

%%%%%%%%%%%%%%%%%%%%%%%%%%%%%%%%%%%%%%%%%%%%%%%%%%%%%%%%%%%%
\subsubsection{Confusing Relative Consistency With Absolute Consistency}

A statement of the form

\[
\operatorname{Con}(T)
\Longrightarrow
\operatorname{Con}(T+\varphi)
\]

does not prove

\[
\operatorname{Con}(T).
\]

It only compares consistency strengths.

%%%%%%%%%%%%%%%%%%%%%%%%%%%%%%%%%%%%%%%%%%%%%%%%%%%%%%%%%%%%
\subsubsection{Thinking the Cumulative Hierarchy Is a Set}

The full hierarchy

\[
V
\]

is a proper class.

Each individual stage

\[
V_\alpha
\]

is a set.

%%%%%%%%%%%%%%%%%%%%%%%%%%%%%%%%%%%%%%%%%%%%%%%%%%%%%%%%%%%%
\subsubsection{Assuming Every First-Order Model of ZFC Looks Like the Intended Universe}

Model theory permits nonstandard models.

Questions about internal and external cardinality, natural numbers, and
well-foundedness require careful distinction.

%%%%%%%%%%%%%%%%%%%%%%%%%%%%%%%%%%%%%%%%%%%%%%%%%%%%%%%%%%%%
\subsection{Applications}
\label{subsec:axiomatic-set-theory-applications}
%%%%%%%%%%%%%%%%%%%%%%%%%%%%%%%%%%%%%%%%%%%%%%%%%%%%%%%%%%%%

\begin{applicationbox}

\textbf{Foundations of Mathematics.}
ZFC provides a common axiomatic framework in which most ordinary
mathematics can be formalized.

\medskip

\textbf{Analysis.}
Real numbers, functions, measures, sigma-algebras, and function spaces can
all be represented set-theoretically.

\medskip

\textbf{Algebra.}
Groups, rings, fields, modules, and vector spaces are sets equipped with
operations satisfying axioms.

\medskip

\textbf{Topology.}
A topology is a set of subsets satisfying structural conditions, making
power sets and Choice central foundational tools.

\medskip

\textbf{Functional Analysis.}
The Axiom of Choice and equivalent principles appear in results involving
bases, maximal families, and extension theorems.

\medskip

\textbf{Model Theory.}
Models of set theory provide the setting for compactness, independence,
and nonstandard phenomena.

\medskip

\textbf{Category Theory.}
Set/class distinctions and universe constructions help control
size-related foundational issues.

\medskip

\textbf{Set Theory Itself.}
Ordinals, cardinals, forcing, constructibility, large cardinals, and
independence questions form a major research area.

\end{applicationbox}

%%%%%%%%%%%%%%%%%%%%%%%%%%%%%%%%%%%%%%%%%%%%%%%%%%%%%%%%%%%%
\subsection{Exercises}
\label{subsec:axiomatic-set-theory-exercises}
%%%%%%%%%%%%%%%%%%%%%%%%%%%%%%%%%%%%%%%%%%%%%%%%%%%%%%%%%%%%

\begin{exercise}[1]
State the Axiom of Extensionality.
\end{exercise}

\begin{exercise}[1]
State the Axiom of Pairing.
\end{exercise}

\begin{exercise}[1]
State the Axiom of Union.
\end{exercise}

\begin{exercise}[1]
State the Axiom of Power Set.
\end{exercise}

\begin{exercise}[1]
Explain the Separation Schema.
\end{exercise}

\begin{exercise}[1]
Explain the Replacement Schema.
\end{exercise}

\begin{exercise}[1]
State the Axiom of Infinity.
\end{exercise}

\begin{exercise}[1]
Explain the Axiom of Foundation.
\end{exercise}

\begin{exercise}[1]
State the Axiom of Choice.
\end{exercise}

\begin{exercise}[1]
Distinguish ZF and ZFC.
\end{exercise}

\begin{exercise}[2]
Explain Russell's paradox.
\end{exercise}

\begin{exercise}[2]
Explain why Russell's paradox rules out unrestricted comprehension.
\end{exercise}

\begin{exercise}[2]
Prove that the empty set is unique using Extensionality.
\end{exercise}

\begin{exercise}[2]
Use Pairing to construct a singleton.
\end{exercise}

\begin{exercise}[2]
Use Pairing and Union to construct \(A\cup B\).
\end{exercise}

\begin{exercise}[2]
Construct

\[
A\cap B
\]

using Separation.
\end{exercise}

\begin{exercise}[2]
Construct

\[
A\setminus B
\]

using Separation.
\end{exercise}

\begin{exercise}[2]
List the elements of

\[
\mathcal{P}(\{a,b,c\}).
\]
\end{exercise}

\begin{exercise}[2]
Explain why

\[
|\mathcal{P}(A)|=2^{|A|}
\]

for finite \(A\).
\end{exercise}

\begin{exercise}[3]
Explain the difference between Separation and Replacement.
\end{exercise}

\begin{exercise}[3]
Give an example of a definable image whose existence is naturally justified
using Replacement.
\end{exercise}

\begin{exercise}[2]
Construct the first five von Neumann natural numbers.
\end{exercise}

\begin{exercise}[2]
Verify that

\[
n=\{0,1,\ldots,n-1\}
\]

for the first several finite ordinals.
\end{exercise}

\begin{exercise}[3]
Show that

\[
m<n
\iff
m\in n
\]

for finite von Neumann ordinals.
\end{exercise}

\begin{exercise}[3]
Explain why the Axiom of Infinity gives at least one infinite set.
\end{exercise}

\begin{exercise}[3]
Explain how \(\omega\) is obtained as the least inductive set.
\end{exercise}

\begin{exercise}[3]
Prove that Foundation implies

\[
x\notin x.
\]
\end{exercise}

\begin{exercise}[3]
Show that Foundation rules out a two-cycle

\[
x\in y\in x.
\]
\end{exercise}

\begin{exercise}[2]
Construct a choice function for a finite family of nonempty sets.
\end{exercise}

\begin{exercise}[3]
Explain why finite choice does not capture the full Axiom of Choice.
\end{exercise}

\begin{exercise}[3]
Explain the difference between a choice function and an arbitrary function.
\end{exercise}

\begin{exercise}[2]
Verify the Kuratowski ordered-pair identity for a simple example.
\end{exercise}

\begin{exercise}[3]
Prove that

\[
(a,b)=(c,d)
\]

implies

\[
a=c
\]

and

\[
b=d.
\]
\end{exercise}

\begin{exercise}[2]
Define the Cartesian product set-theoretically.
\end{exercise}

\begin{exercise}[2]
Define a binary relation as a set.
\end{exercise}

\begin{exercise}[2]
Define a function as a set of ordered pairs.
\end{exercise}

\begin{exercise}[3]
Explain why functions need no new primitive ontology beyond sets.
\end{exercise}

\begin{exercise}[2]
Determine whether a given set is transitive.
\end{exercise}

\begin{exercise}[3]
Show that every von Neumann natural number is transitive.
\end{exercise}

\begin{exercise}[3]
Explain why an ordinal is the set of all smaller ordinals.
\end{exercise}

\begin{exercise}[2]
Compute the successor of a given finite ordinal.
\end{exercise}

\begin{exercise}[3]
Explain why \(\omega\) is a limit ordinal.
\end{exercise}

\begin{exercise}[3]
Show conceptually that

\[
1+\omega=\omega.
\]
\end{exercise}

\begin{exercise}[3]
Explain why

\[
\omega+1\neq\omega.
\]
\end{exercise}

\begin{exercise}[3]
State transfinite induction.
\end{exercise}

\begin{exercise}[3]
Explain the base, successor, and limit stages of transfinite induction.
\end{exercise}

\begin{exercise}[3]
Explain transfinite recursion conceptually.
\end{exercise}

\begin{exercise}[2]
Define equipotence.
\end{exercise}

\begin{exercise}[2]
Construct a bijection between \(\mathbb{N}\) and the even natural numbers.
\end{exercise}

\begin{exercise}[3]
Explain why a proper subset of an infinite set may have the same
cardinality as the entire set.
\end{exercise}

\begin{exercise}[3]
Construct a bijection between \(\mathbb{N}\) and \(\mathbb{Z}\).
\end{exercise}

\begin{exercise}[3]
Explain why \(\mathbb{Q}\) is countable.
\end{exercise}

\begin{exercise}[3]
State Cantor's theorem.
\end{exercise}

\begin{exercise}[3]
Prove Cantor's theorem using the diagonal set

\[
D
=
\{x\in A:x\notin f(x)\}.
\]
\end{exercise}

\begin{exercise}[3]
Use Cantor's theorem to show there is no set of all sets.
\end{exercise}

\begin{exercise}[2]
Define the continuum cardinal \(\mathfrak{c}\).
\end{exercise}

\begin{exercise}[3]
Explain why

\[
\mathfrak{c}
=
2^{\aleph_0}.
\]
\end{exercise}

\begin{exercise}[2]
State the Continuum Hypothesis.
\end{exercise}

\begin{exercise}[3]
State the Generalized Continuum Hypothesis.
\end{exercise}

\begin{exercise}[3]
Explain cardinal addition using disjoint unions.
\end{exercise}

\begin{exercise}[3]
Explain cardinal multiplication using Cartesian products.
\end{exercise}

\begin{exercise}[3]
Explain cardinal exponentiation using function sets.
\end{exercise}

\begin{exercise}[2]
Define a well-order.
\end{exercise}

\begin{exercise}[3]
State the Well-Ordering Theorem.
\end{exercise}

\begin{exercise}[3]
State Zorn's Lemma.
\end{exercise}

\begin{exercise}[3]
Distinguish maximal and maximum elements.
\end{exercise}

\begin{exercise}[3]
Give a poset with several maximal elements and no maximum.
\end{exercise}

\begin{exercise}[3]
Outline the standard Zorn's Lemma proof strategy.
\end{exercise}

\begin{exercise}[3]
Explain why the union of a chain of linearly independent sets is linearly
independent.
\end{exercise}

\begin{exercise}[3]
Use that fact to explain why Zorn's Lemma implies every vector space has a
basis.
\end{exercise}

\begin{exercise}[3]
State the equivalence of AC, Zorn's Lemma, and the Well-Ordering Theorem.
\end{exercise}

\begin{exercise}[3]
Define a filter.
\end{exercise}

\begin{exercise}[3]
Define an ultrafilter.
\end{exercise}

\begin{exercise}[3]
Explain why ultrafilters are maximal proper filters.
\end{exercise}

\begin{exercise}[2]
Define the transitive closure of a set conceptually.
\end{exercise}

\begin{exercise}[3]
Compute the first few stages of the transitive closure of a simple finite
set.
\end{exercise}

\begin{exercise}[2]
Compute

\[
V_0,
\qquad
V_1,
\qquad
V_2.
\]
\end{exercise}

\begin{exercise}[3]
Compute \(V_3\).
\end{exercise}

\begin{exercise}[3]
Explain the successor and limit stages of the cumulative hierarchy.
\end{exercise}

\begin{exercise}[3]
Explain why every member of \(V_{\alpha+1}\) is a subset of \(V_\alpha\).
\end{exercise}

\begin{exercise}[3]
Explain why membership decreases rank.
\end{exercise}

\begin{exercise}[3]
Explain why \(V\) is a proper class rather than a set.
\end{exercise}

\begin{exercise}[3]
Explain why the collection of all ordinals is a proper class.
\end{exercise}

\begin{exercise}[3]
Describe the Burali--Forti paradox.
\end{exercise}

\begin{exercise}[3]
Explain what a model of ZFC is.
\end{exercise}

\begin{exercise}[3]
Distinguish internal and external statements about a model.
\end{exercise}

\begin{exercise}[3]
Define relative consistency.
\end{exercise}

\begin{exercise}[3]
Define independence from a theory.
\end{exercise}

\begin{exercise}[3]
Explain why

\[
T\nvdash\varphi
\]

alone does not prove that \(\varphi\) is independent of \(T\).
\end{exercise}

\begin{exercise}[3]
State the independence status of CH relative to ZFC.
\end{exercise}

\begin{exercise}[3]
Explain the constructible universe \(L\) conceptually.
\end{exercise}

\begin{exercise}[3]
Explain forcing conceptually without technical construction.
\end{exercise}

\begin{exercise}[3]
Explain why forcing is useful for independence proofs.
\end{exercise}

\begin{exercise}[3]
Explain what a large-cardinal axiom is conceptually.
\end{exercise}

\begin{exercise}[3]
Explain why large cardinals are viewed as stronger existence principles.
\end{exercise}

\begin{exercise}[4]
Study the derivation of the empty set from other standard ZF axioms.
\end{exercise}

\begin{exercise}[4]
Formalize each ZF axiom in first-order logic.
\end{exercise}

\begin{exercise}[4]
Study the precise logical form of the Separation Schema.
\end{exercise}

\begin{exercise}[4]
Study the precise logical form of the Replacement Schema.
\end{exercise}

\begin{exercise}[4]
Prove the recursion theorem on \(\omega\).
\end{exercise}

\begin{exercise}[4]
Develop ordinal addition recursively.
\end{exercise}

\begin{exercise}[4]
Develop ordinal multiplication recursively.
\end{exercise}

\begin{exercise}[4]
Develop ordinal exponentiation recursively.
\end{exercise}

\begin{exercise}[4]
Prove transfinite induction.
\end{exercise}

\begin{exercise}[4]
Study the transfinite recursion theorem.
\end{exercise}

\begin{exercise}[4]
Study cardinal numbers as initial ordinals.
\end{exercise}

\begin{exercise}[4]
Prove the Cantor--Bernstein theorem.
\end{exercise}

\begin{exercise}[4]
Study cardinal arithmetic with and without Choice.
\end{exercise}

\begin{exercise}[4]
Prove equivalence of the Axiom of Choice and the Well-Ordering Theorem.
\end{exercise}

\begin{exercise}[4]
Prove equivalence of Zorn's Lemma and the Axiom of Choice.
\end{exercise}

\begin{exercise}[4]
Study the Ultrafilter Lemma and compare its strength with full Choice.
\end{exercise}

\begin{exercise}[4]
Develop the rank function rigorously.
\end{exercise}

\begin{exercise}[4]
Prove that every set belongs to some \(V_\alpha\).
\end{exercise}

\begin{exercise}[4]
Study transitive models of set theory.
\end{exercise}

\begin{exercise}[4]
Study Gödel's constructible universe \(L\).
\end{exercise}

\begin{exercise}[4]
Study Cohen forcing conceptually.
\end{exercise}

\begin{exercise}[4]
Study the independence of CH.
\end{exercise}

\begin{exercise}[4]
Study inaccessible cardinals.
\end{exercise}

\begin{exercise}[4]
Study measurable cardinals.
\end{exercise}

\begin{exercise}[4]
Investigate alternative set theories with classes.
\end{exercise}

%%%%%%%%%%%%%%%%%%%%%%%%%%%%%%%%%%%%%%%%%%%%%%%%%%%%%%%%%%%%
\subsection{Conceptual Summary}
\label{subsec:axiomatic-set-theory-summary}
%%%%%%%%%%%%%%%%%%%%%%%%%%%%%%%%%%%%%%%%%%%%%%%%%%%%%%%%%%%%

Axiomatic set theory replaces unrestricted collection formation with a
controlled formal theory.

The basic language uses membership:

\[
\boxed{
\in.
}
\]

The central ZF axioms govern:

\[
\boxed{
\text{identity of sets}
+
\text{set construction}
+
\text{infinity}
+
\text{well-foundedness}.
}
\]

Separation replaces unrestricted comprehension with

\[
\boxed{
\{x\in A:\varphi(x)\}.
}
\]

Replacement permits definable images of sets.

Infinity provides the natural-number hierarchy:

\[
\boxed{
0=\varnothing,
\qquad
n+1=n\cup\{n\}.
}
\]

Ordinals extend counting and ordering into the transfinite.

Cardinals measure size up to bijection.

Cantor's theorem gives

\[
\boxed{
|A|
<
|\mathcal{P}(A)|
}
\]

for every set \(A\), generating an unending hierarchy of infinite
cardinalities.

Adding the Axiom of Choice produces

\[
\boxed{
\mathrm{ZFC}.
}
\]

Choice is equivalent over ZF to major principles such as

\[
\boxed{
\text{Zorn's Lemma}
}
\]

and

\[
\boxed{
\text{the Well-Ordering Theorem}.
}
\]

The cumulative hierarchy

\[
\boxed{
V_0=\varnothing,
\qquad
V_{\alpha+1}=\mathcal{P}(V_\alpha),
\qquad
V_\lambda=\bigcup_{\alpha<\lambda}V_\alpha
}
\]

organizes the set-theoretic universe by rank.

Advanced set theory then studies

\[
\boxed{
\text{models}
+
\text{independence}
+
\text{forcing}
+
\text{constructibility}
+
\text{large cardinals}.
}
\]

%%%%%%%%%%%%%%%%%%%%%%%%%%%%%%%%%%%%%%%%%%%%%%%%%%%%%%%%%%%%
\subsection{Section Connections}
\label{subsec:axiomatic-set-theory-connections}
%%%%%%%%%%%%%%%%%%%%%%%%%%%%%%%%%%%%%%%%%%%%%%%%%%%%%%%%%%%%

\chapterconnection{
Axiomatic Set Theory builds directly on the formal-language framework of
Section~\ref{sec:formal-logic}. ZF and ZFC are first-order theories whose
single primitive nonlogical relation is membership. The axioms justify the
set constructions used throughout mathematics, while ordinals, cardinals,
Choice, the cumulative hierarchy, and independence phenomena reveal the
deeper structure of the mathematical universe. The next section, Model
Theory, studies formal theories through their mathematical models and
develops tools for comparing structures, elementary equivalence,
definability, compactness, and model construction. Later sections on Proof
Theory and Computability Theory study the proof-theoretic and algorithmic
sides of the same foundational framework.
}

%%%%%%%%%%%%%%%%%%%%%%%%%%%%%%%%%%%%%%%%%%%%%%%%%%%%%%%%%%%%
% SECTION SOURCE TRACKING
%%%%%%%%%%%%%%%%%%%%%%%%%%%%%%%%%%%%%%%%%%%%%%%%%%%%%%%%%%%%

%------------------------------------------------------------
% 12.2 Axiomatic Set Theory
%------------------------------------------------------------
%
% Source basis:
%   - Standard Zermelo--Fraenkel set theory.
%   - Standard set-theoretic treatment of ordinals, cardinals, Choice,
%     cumulative hierarchy, and independence.
%
% Used for:
%   - naive versus axiomatic set theory;
%   - Russell's paradox;
%   - language of set theory;
%   - Extensionality;
%   - empty set;
%   - Pairing;
%   - Union;
%   - Power Set;
%   - Separation;
%   - Replacement;
%   - Infinity;
%   - Foundation;
%   - Axiom of Choice;
%   - ZF and ZFC;
%   - derived set constructions;
%   - ordered pairs;
%   - Cartesian products;
%   - relations and functions as sets;
%   - von Neumann natural numbers;
%   - inductive sets and omega;
%   - transitive sets;
%   - ordinals;
%   - successors and limits;
%   - ordinal ordering and arithmetic;
%   - transfinite induction;
%   - transfinite recursion;
%   - equipotence;
%   - cardinal numbers;
%   - countability;
%   - Cantor's theorem;
%   - no universal set;
%   - continuum cardinality;
%   - aleph numbers;
%   - Continuum Hypothesis;
%   - Generalized Continuum Hypothesis;
%   - cardinal arithmetic;
%   - well-orders;
%   - Well-Ordering Theorem;
%   - Zorn's Lemma;
%   - maximal versus maximum elements;
%   - vector-space basis application;
%   - equivalent forms of Choice;
%   - filters and ultrafilters;
%   - transitive closure;
%   - rank;
%   - cumulative hierarchy;
%   - proper classes;
%   - class of all ordinals;
%   - Burali--Forti paradox;
%   - models of set theory;
%   - internal and external perspectives;
%   - relative consistency;
%   - independence;
%   - CH independence;
%   - constructible universe preview;
%   - forcing preview;
%   - large-cardinal preview;
%   - applications and exercises.
%
% Notes:
%   - Model Theory is developed next in Section 12.3.
%   - Formal Logic appears in Section 12.1.
%   - Proof Theory follows in Section 12.4.
%   - Computability Theory follows in Section 12.5.
%
%%%%%%%%%%%%%%%%%%%%%%%%%%%%%%%%%%%%%%%%%%%%%%%%%%%%%%%%%%%%